% =====================================================================
%  Harald Høffding, BEMÆRKNINGER OM DEN PLATONISKE DIALOG PARMENIDES.
%  Det Kgl. Danske Videnskabernes Selskab, Filosofiske Meddelelser I, 2.  København: Hovedkommissionær Andr. Fred.
%  Høst & Søn, Kgl. Hof-Boghandel; Bianco Lunos Bogtrykkeri, 1920.  65 pp.  Antiqua.
%
%  TRANSCRIPTION -- GENERATED; DO NOT EDIT BY HAND.  Rebuild from the repo root:
%    sh ebook-method/platons-parmenides/rebuild.sh   (that is:)
%    python3 ebook-method/platons-parmenides/draft.py --emit --force \
%      && python3 ebook-method/collate/check.py platons-parmenides --italics --apply \
%      && python3 ebook-method/collate/check.py platons-parmenides --apply \
%      && python3 ebook-method/platons-parmenides/draft.py --post
%
%  STATUS: COMPLETE (2026-10-20): title page, the whole text pp. 3-65 with its 34 notes, and the print's Indhold.
%
%  EDITION.  The text of record is the Royal Danish Library's scan (~/bibliotek/Høffding, Harald/1920-platoniske-dialog-
%  parmenides.pdf, SCANS.tsv by sha256), 400 ppi colour, with an ABBYY FineReader Server text layer.  The only printing; no
%  e-book, no reprint and no earlier transcription found (checked 2026-10-20: the repo, Internet Archive, Google Books, NB,
%  the web).  (Mindre Arbejder II, 1905, holds a different piece, »Parmenides fra Elea«, 1901.)
%
%  ACCURACY (2026-10-20).  The words are those of the scan's layer (draft.py), collated word by word with an independent
%  tesseract reading of the images (collate/check.py, witness 'ocr'): 384 disagreements in ~17,350 words at the first
%  collation.  check.py's own rules settled 35; the Greek (177 words, which the layer garbles into Latin letters) was set
%  from tesseract's grc model at 400 dpi (greek.tsv) and settled by rule (136 checks, every non-Greek word of the check
%  required in both readings); rules.py's candidate rules were switched on only after a BLIND PILOT (reader A, 52 crops,
%  12/12 controls) measured them on this book: the layer's glyph faults (l for t/f, li for h) 8/8, tesseract's 8/8, small
%  capitals 8/8, note figures 5/5, the note-seam rule narrowed after a miss.  The rest were read at the image BLIND
%  (sheets.py: 129 crops, 26/26 controls).  OPEN 0 on a fresh draft.
%  STRUCTURE read at the image: every note whole (reader S, n1 sheets; pp. 22, 32, 36, 37, 50 in the main session) and set
%  from that reading (patches.NOTES), the marks placed by the draft and checked page by page against the notes; the Greek
%  read blind at 600 dpi (62 crops) and every doubtful breathing or accent at 800 dpi in the main session; the title page,
%  the Indhold, the four chapter heads and the large initial P.  SWEEPS after OPEN 0, each spot read at the image:
%  check.py --moved (a line put a line up, p. 9), --splits, --dropped (note seams only); the unknown-word sweep (the
%  misprints below; „2111)“ for „211 D“ inside a Greek check, p. 47); the quote-balance sweep (two notes the draft had
%  left in the text, pp. 32 and 50); small capitals (33 crops: names are in small capitals at their first mention and in
%  the notes only); punctuation after note figures (three stops and commas the layer lost, pp. 17, 22, 25, found by the
%  translators).  SAMPLED PAGE CHECK (pagecheck.py, seed 202610201): pp. 16, 34, 56 and the title page read whole and
%  blind, 889 words: 0 word errors (95 % upper bound 3.4 per 1000), 1 punctuation error (a speck read as an apostrophe by
%  both OCRs, p. 34, mended).  Page markers: 63 of 63 exact (check.py --pages).  The ø rule (Hans 2026-10-08): nothing to apply.
%  What this cannot see: a misprint both OCRs normalise into a real word.
%  PRINTED AS IS (each seen at the image; a % comment at the site): „Invendinger“ (p. 4), „Aristotetes’“ (p. 17),
%  „for-|forsøger“ (p. 19), „«Staten«“ with a turned opening guillemet (p. 32), „BURNÉT“ (p. 62, note).  The Greek as
%  printed, accents and breathings included: αὐτὰς, αὐτά (p. 23, the editions αὑτ-), ὤν (pp. 23, 59), ἥ for ἡ (pp. 26, 50),
%  τελευτᾲ (p. 24), αὖξηθείς (p. 46), ἐίπερ (p. 47), οὖδὲ (p. 50); the print's curly theta ϑ is set θ (textalpha has no
%  text-mode ϑ).  ἐστιν (p. 23, note) and ἐξαίφνης (p. 39) carry a ^-shaped mark over the ε in the print: set with the
%  smooth breathing.  Edition figures as printed: »Søren Kierkegaard som Filosof«² (p. 25), Vorsokratiker’ (p. 57),
%  Platonische Studien? (p. 58).  PRINTER'S DEFECT: „indgaaende“ (p. 54), its second e broken.
%  The Indhold gives the Indledning as p. 1 (the text begins on p. 3) and chapter 3 as „Begrebet det Pludselige i Platon’s
%  Filosofi“ (the head, p. 38: „… (τὸ ἐξαίφνης) i Platons Filosofi.“); chapter 4's head quotes „Parmenides“ with „ “, the
%  text with » «.  Both readings recorded as printed.
%  Not transcribed: running heads, signatures, the series line „Vidensk. Selsk. Filosof. Medd. I, 2.“, the Society's seal
%  (described).  The large initial P (p. 3) is set as a plain letter, „LATON“ after it in small capitals as printed.
%
%  PAGE MAP.  printed p = PDF page - 2: title p. [1] = PDF 3, text pp. 3-65 = PDF 5-67, Indhold = PDF 68.  Chapter heads:
%  1. p. 6, 2. p. 26, 3. p. 38, 4. p. 53.
%
%  TOOLS (ebook-method/platons-parmenides/, copied from om-samfundets-vilkaar): draft.py, ocr.py, greek.tsv + greek_ocr.py,
%  rules.py, sheets.py + reads.py + READER-BRIEF.md, structsheets.py, patches.py (PATCHES, NOTES, POST), decisions.tsv,
%  pagecheck.py + PAGECHECK-BRIEF.md, rebuild.sh; the translation: GLOSSARY.md, TRANSLATOR-BRIEF.md, tr_extract.py,
%  tr_build.py, tr_template.tex, tr/.
%
% =====================================================================
\documentclass[12pt,a4paper,openany]{book}
\usepackage[utf8]{inputenc}
\usepackage[T1]{fontenc}
\usepackage{libertinus}
\usepackage{libertinust1math}
\usepackage{textalpha}
\usepackage[danish]{babel}
\usepackage[protrusion=true,expansion=true]{microtype}
\usepackage{setspace}
\usepackage[margin=1.3in]{geometry}
\usepackage{fancyhdr}
\usepackage{hyperref}
\hypersetup{pdftitle={Bemærkninger om den platoniske Dialog Parmenides}, pdfauthor={Harald Høffding}, hidelinks}
\pagestyle{fancy}
\fancyhf{}
\fancyhead[LE,RO]{\thepage}
\fancyhead[C]{\textit{Høffding: Bemærkninger om den platoniske Dialog Parmenides (1920)}}
\setstretch{1.3}
\usepackage{marginnote}
\newcommand{\tocline}[3]{\noindent\makebox[2.2em][r]{#1}\hspace{0.6em}#2\dotfill\ #3\par}
\newcommand{\opage}[1]{\marginnote{\footnotesize\textit{[#1]}}}
% a chapter heading, centred and bold, as printed (the print's face is a bold grotesque); the large initial P (p. 3) is not reproduced.
\newcommand{\chaphead}[1]{\par\bigskip\phantomsection\begin{center}\bfseries #1\end{center}\medskip\ignorespaces}

\begin{document}

% --- title page (PDF 3, p. [1]), read at the image (m1 F3, reader S, 2026-10-20).  PDF 1, the series cover, prints the same
% title page; PDF 2 the series notice; PDF 4 (p. [2]) is blank.
\begin{titlepage}
\begin{center}
Det Kgl. Danske Videnskabernes Selskab.\\[0.2em]
{\small Filosofiske Meddelelser \textbf{I}, 2.}\\[0.3em]
\rule{\textwidth}{0.4pt}\\[5em]
{\LARGE BEMÆRKNINGER}\\[0.4em]
{\LARGE OM DEN PLATONISKE DIALOG}\\[0.4em]
{\LARGE PARMENIDES}\\[2em]
{\small AF}\\[2em]
{\large HARALD HØFFDING}\\[4em]
[the Society's seal]\\[5em]
{\Large KØBENHAVN}\\[0.6em]
{\small\textsc{Hovedkommissionær:}} ANDR. FRED. HØST \& SØN, {\small\textsc{Kgl. Hof-Boghandel}}\\[1.2em]
BIANCO LUNOS BOGTRYKKERI\\[0.3em]
1920
\end{center}
\end{titlepage}
\clearpage


% ===== TEXT (pp. 3--65) =====
% --- p. 3 ---
\setcounter{footnote}{0}%
\opage{3}P\textsc{laton}’s »Parmenides« er en af de mest omstridte af
hans Dialoger. Den falder i to Afsnit, der hvert for
sig frembyder Vanskeligheder for Platonforskningen.

I det første Afsnit udvikler Parmenides fra Elea overfor
Sokrates, der her fremstilles som en ung Mand, forskellige
Indvendinger, som synes at ramme den Filosofi,
den saakaldte Idelære, som Platon har fremsat i sine berømteste
Dialoger, dem, som alle er enige om at betragte som
Højdepunktet i hans Forfattervirksomhed: »Symposion«,
»Faidon« og »Staten«. Hvis nu disse Indvendinger virkelig
rammer, hvilket synes at blive indrømmet af Platon
selv, idet han lader Sokrates staa uden Udveje med Hensyn
til sin fortsatte Filosoferen, hvorledes kan det da forklares,
at Platon har fremsat dem? Dette Spørgsmaal er
blevet besvaret paa forskellig Maade. En hel Række tyske
Forskere (\textsc{Schaarschmidt}, \textsc{Überweg} o. fl.) antog, at »Parmenides«
slet ikke var af Platon, fordi Indholdet var uforeneligt
med Indholdet af de almindelig anerkendte platoniske
Dialoger. Andre (\textsc{Zeller}, \textsc{Jackson}, \textsc{Natorp}, \textsc{Ræder})
søgte at vise, at Platon selv, dels i Dialogens andet Afsnit,
dels i senere Dialoger, har besvaret de fremførte Indvendinger;
men i Svarets Indhold afviger saa disse Forskere
indbyrdes. Andre endelig (\textsc{Apelt}, og nylig \textsc{Wilamowitz-Moellendorff}
 i sin Platonbiografi) betragter Indvendingerne
% --- p. 4 ---
\setcounter{footnote}{0}%
\opage{4}som Øvelseseksempler, Opgaver for Platon’s Disciple, idet
Platon har stolet paa, at de for de kløgtigste af disse vilde
falde bort af sig selv.

For mit Vedkommende maa jeg efter gentagen Beskæftigelse
med det mærkelige Skrift være enig med dem, der
% PRINTED AS IS: „Invendinger“ (for Indvendinger), p. 4, seen at the image
mener, at de Invendinger, der gøres i første Afsnit virkelig
rammer. Men jeg mener, man maa gaa et Skridt videre.
Platon har i disse Indvendinger --- hvad enten de stammer
fra andre eller fra ham selv --- staaet overfor Synsmaader,
som ikke blot har Interesse ved Kritik af hans egen Filosofi,
men er af Interesse overfor ethvert Forsøg paa Metafysik
(i Betydning af afsluttende Erkendelse af Tilværelsen). Jeg
vil søge at godtgøre dette gennem Analyse af Dialogens
første Del. Jeg bemærker straks, at Platon herved er vokset
i mine Øjne. Vel har han ikke været og ej heller kunnet
være sig Rækkevidden af de fremførte Indvendinger bevidst,
og de har heller ikke ført ham til helt at opgive Idelæren.
Men indenfor den mærkelige Blanding af Kritik og Mystik,
som udgør Platon’s Ejendommelighed, er i »Parmenides«
det kritiske Element momentvis kommet frem med saa
stor Styrke, at han senere ikke har formaaet at bøje det
ind som harmonisk Led i sit System, saa betydningsfulde
end i flere Henseender de nye Undersøgelser er, han i
senere Dialoger slaar ind paa. Man forstaar, hvorledes der
af Platon’s Skole kunde fremgaa en kritisk Tænker som
Karneades, Oldtidens Hume, medens Platon gennem det
andet Element i sin Filosofi blev Stamfader til nyplatonisk
og middelalderlig Mystik.

Dialogens andet Afsnit vil give en Anvisning til fortsat
Filosoferen, selv om man indrømmer de i første
Afsnit fremførte Indvendinger mod Idelæren. Den Vej, der
følges, er den hypotetiske Slutnings. Der undersøges forskel%
% --- p. 5 ---
\setcounter{footnote}{0}%
\opage{5}lige, indbyrdes stridende Forudsætninger; disses Konsekvenser
skal blive afgørende for, hvilke Synsmaader man
holder sig til. Men Fremstillingen er i sine fleste Dele saa
uklar, tillige saa abstrakt, og hvad der er værre, saa sofistisk,
at man har haft Vanskelighed ved at tillægge Platon
den eller ved at anse den for alvorlig ment. De, der forkaster
Dialogens Ægthed, henter mægtige Vaaben fra dette
Afsnit. De, der allerede i første Afsnit blot saa et Øvelseseksempel,
mener, at dette i endnu højere Grad gælder om
dette andet Afsnit. Men der er dog ogsaa Forskere, som
her har fundet Momenter til virkelig Besvarelse af det
første Afsnits Indvendinger, i hvert Tilfælde Udveje til at
komme udenom dem.

For mit eget Vedkommende tilstaar jeg, at jeg er ude
af Stand til at finde nogen Betydning i største Delen af
de meget abstrakte og tildels paa afgjorte Fejlslutninger
byggede Ræsonnementer, der udgør Dialogens andet Afsnit.
Men i Løbet af den uklare Tankegang fremtræder der mærkværdige
Tankeglimt, som bør fremdrages, og som tildels
faar Interesse ved at sammenholdes med andre Tanker
hos Platon. Som Helhed gør Afsnittet et meget ufærdigt
Indtryk. Der mangler baade formal og real Afslutning, og
det er ikke usandsynligt, at der her egentlig kun foreligger
et Udkast eller et Skema, der kunde berigtiges og tydeliggøres
ved virkelig Udførelse. Maaske har ivrige Disciple
ubetænksomt udgivet et af Platon efterladt Udkast. ---

Den saaledes antydede Opfattelse af den mærkelige
Dialogs to Dele forsøger jeg nu ved nærmere Gennemgang
at begrunde.

% --- p. 6 ---
\setcounter{footnote}{0}%
\opage{6}\chaphead{1. Kritik af Idelæren.}


a. Ved en Sammenkomst i Athen mellem Eleaterne
Parmenides og Zenon paa den ene Side, Sokrates (som
Platoniker) paa den anden Side, --- en Sammenkomst, som
det er meget vanskeligt at gøre historisk mulig, --- oplæses
først \textsc{Zenon}’s berømte (desværre nu tabte) Skrift, i hvilket
han søger at godtgøre Modsigelserne ved den Antagelse, at
Tilværelsen skulde bestaa af Forskelligheder, altsaa være
en Mangfoldighed. De berømte Indvendinger mod Bevægelsens
Mulighed, som især har vakt Opsigt, er kun Følgeslutninger
af denne Kritik: Bevægelse forudsætter Mangfoldighed
i Rum og Tid, og Vanskeligheden er da, hvorledes
man kan komme over fra en Rums- og Tidsdel til
en anden. Efter Oplæsningen siger Sokrates, at Zenon’s
Skrift indeholder det samme som Parmenides’ Læredigt;
Forskellen er kun, at Parmenides ligefrem erklærer, at
kun det Ene, det med sig selv lige og absolut Ensartede,
eksisterer, og at Flerhed og Forandring er Illusioner,
medens Zenon gaar indirekte frem ved at vise, at Antagelsen
af Flerhed og Vorden fører til Modsigelser. Dette indrømmer
Zenon, og Sokrates hævder saa, at de Modsigelser,
Parmenides’ Modstandere finder i Enhedslæren, og de
Modsigelser, Zenon finder i Flerhedslæren, falder bort,
naar man antager, at Enhed og Flerhed, Lighed og Forskel,
Hvile og Bevægelse er ligesaa mange Ideer, ligesaa mange
% --- p. 7 ---
\setcounter{footnote}{0}%
\opage{7}Arter af Væren, og at de enkelte Emner (Erfaringsgenstande)
hvert for sig i visse Henseender hører ind under
en af disse Arter, i visse andre Henseender under en anden,
(p. 127 E—129 E).

I denne Tankegang fremtræder Modsætningen mellem
Eleatisme og Platonisme klart. Ifølge Parmenides skulde
Tilværelsen kunne tænkes og være tænkt ved en eneste
Tanke, Enhedens, Lighedens (Identitetens) og Ubevægelighedens
Tanke. Platon’s Idelære betegner her et Fremskridt,
idet den hævder, at Tilværelsen kun forstaas ved en Flerhed
af Tanker, hvis indbyrdes Forhold det gælder at bestemme.
Idelæren er bleven til ved Trangen til at danne
Begreber for hver enkelt afgrænset Gruppe af Emner, altsaa
ved en Opstigen fra en Flerhed af Emner til en Enhed
i Tanken, en Opstigen, der maatte gentages, hver
Gang man stod overfor en ny Gruppe af Emner. Og Platon
lader nu Sokrates hævde, at hvert Emne eller hver Emnegruppe
kan falde ind under forskellige Tanker eller Ideer,
idet der ingen Modsigelse er i, at forskellige Ideer mødes
i samme Emne. De eleatiske Kritikere gaar ikke nærmere
ind paa denne Side af Sagen, saa vigtig den i og for sig
er. Hvad Platon her gør gældende, kan betegnes ved, hvad
H. C. Ørsted har kaldt Samfoldighed, hvorved han forstod
den Maade, hvorpaa forskellige Naturlove i deres Sammenspil
udgør det enkelte Emnes (Fænomens) Væsen. Springvandet
f. Eks. bliver os forstaaeligt ved Foreningen af Lovene
for Fald, Draabedannelse, Lys og Lyd o. s. v. Ørsted
bruger ogsaa Udtrykket Ide om en saadan Samfoldighed.

Den Kritik, som Parmenides nu i det følgende retter
mod Platonismen, angaar derimod et andet Hovedpunkt,
nemlig Forholdet mellem Ideer og Emner overhovedet.
Platon havde hævdet, at Tilværelsen ikke kan tænkes blot
% --- p. 8 ---
\setcounter{footnote}{0}%
\opage{8}ved Hjælp af det eleatiske Ene. Nu undersøger Parmenides,
om Tilværelsen da kan tænkes ved Hjælp af de platoniske
Ideer.

Først spørger han, om Platon tror, at der gives anden
Lighed end den, vi indenfor vor Erfaring forefinder mellem
forskellige Emner. Det siger Platon ja til. --- Han har jo,
kan vi indskyde, i »Faidon« (p. 74) netop benyttet Lighedsbegrebet
til at vise, hvorledes Erfaringens Emner kan frembyde
forskellige Grader af en og samme Egenskab, og hvorledes
vi ved at følge den opadstigende Graderække, her
altsaa af mere og mere fuldstændige Ligheder, kan komme
til det fuldkomne, absolute Begreb, her altsaa Begrebet
Identitet, som et Ideal, hvortil Erfaringen kun kan frembyde
Trin eller Tilnærmelser.

Fra rent logiske Begreber, paa hvilke Lighed er et Eksempel,
gaar Parmenides over til etiske Begreber, om det
Skønne, det Gode og det Retfærdige. Ogsaa her siger Platon
frejdig ja til Spørgsmaalet, om der gives noget i sig
Skønt, Godt o. s. v., saa at det ikke blot er Erfaringen,
der viser os saadanne Egenskaber. --- Han havde jo ogsaa
i »Faidon« (p. 75 C) udtrykkelig sagt, at saadanne Begreber
som det Skønne, det Gode, det Retfærdige maa betragtes
paa samme Maade som Lighedsbegrebet. Erfaringen
viser os ogsaa her kun ufuldkomne Eksempler, men ved
at fuldende Skalaen naar vi til absolute Ideer.

Anderledes staar Platon (den platoniske Sokrates), naar
Parmenides spørger, om der ogsaa svarer Ideer til Emner
af konkret eller stoflig Natur, som Menneske, Ild, Vand,
Haar, Smuds. Her tør han ikke straks svare ja, skønt han
ser, at det ogsaa her drejer sig om Egenskaber ved foreliggende
Emner, Egenskaber, der kan fremtræde paa mere
eller mindre fuldkommen Maade, saa at samme Tankegang
% --- p. 9 ---
\setcounter{footnote}{0}%
\opage{9}her maa finde Anvendelse som ved de logiske og etiske
Begreber. Han staar tvivlende paa dette Punkt, hvilket
Parmenides erklærer for ungdommelig Usikkerhed (p. 130
B—E).

Dette lille Stykke har stor Interesse ved at kaste Lys
over Idelærens Genesis. Skarpere end andre Steder (smlgn.
især det anførte Sted i »Faidon«) skelnes der her mellem
forskellige Klasser af Ideer. Først stilles de logiske, derpaa
de etiske Grundbegreber (Kategorier) op, og tilsidst de
Artsbegreber, der svarer til Grupper af Enkeltemner. Vi har
vel neppe her en Antydning af den Orden, i hvilken de
forskellige Klasser af Ideer er gaaet op for Platon. Thi det
er sandsynligt (hvad der ogsaa fremgaar af de Dialoger,
der i Almindelighed betragtes som de første), at Platon
som Sokrates’ Discipel først har faaet Øje paa etiske Ideers
Betydning som ledende Tanker, der maatte have evig og
uforanderlig Betydning. Senere er de logiske (erkendelsesteoretiske)
Ideers Betydning bleven ham klar. I »Faidon«
(p. 65 D) begynder han da ogsaa med de etiske Ideer og
gaar derfra over til de logiske (Ligheds- og Identitetsbegrebet),
hvorefter det ganske kort siges, at hvad der gælder
for saadanne Ideer, ogsaa gælder for Begreber som Størrelse,
Sundhed, Styrke, »med ét Ord, for alle Tings Væsen«.
(Lignende Ytring p. 75 C). Det er nu nogle Eksempler paa
disse »alle Ting«, Parmenides holder Platon til Ilden med.
Sagen synes at være den, at Platon ikke fra først af har
skelnet mellem saadanne Begreber, som vi danner om Indholdet
af vore egne intellektuelle Funktioner (Lighed, Størrelse),
og saadanne, som gør sig gældende ved vore etiske
Vurderinger, paa den ene Side, og paa den anden Side
saadanne, som tjener til at ordne og inddele foreliggende
Emner (Genstande, Enkeltvæsener), altsaa Artsbegre%
% --- p. 10 ---
\setcounter{footnote}{0}%
\opage{10}ber i snevrere Betydning. Alle sine Begreber har han vundet
ved Analyse og idealiserende Abstraktion; men i de
Dialoger, der ligger forud for »Parmenides«, var han endnu
ikke opmærksom paa den forskellige Karakter af de saaledes
vundne Begreber. Først senere (i »Theaitetos« og i
»Sofisten«) begynder han at udskille en Gruppe af Begreber,
der svarer til, hvad man kan kalde fundamentale og
formale Kategorier, idet de er umiddelbart knyttede til
selve Tankevirksomheden og derfor finder Anvendelse overalt,
hvor vi tænker, ligegyldigt, hvad vi saa tænker over.
Han lagde herved Grunden til Kategorilæren\footnote[1]{Se om Platon’s Betydning for Kategorilæren \textit{Den menneskelige Tanke,} p. 140 f.}. I Modsætning
til Begreber, vi anvender ved al Tænken, og ligesaa
i Modsætning til dem, vi anvender ved Vurdering af Handlinger
og Mennesker, maatte saa mere og mere de egentlige
Artsbegreber træde, og vi ser da ogsaa, at Platon i
Dialoger fra den følgende Tid (»Sofisten« og »Statsmanden«)
lægger stor Vægt paa Inddeling og Klassifikation
af Emner.

I sit grundige og for mig meget lærerige Skrift »Platon’s
philosophischer Entwickelungsgang« mener Dr. \textsc{Hans}
\textsc{Ræder}, at Idelæren mister en Del af sin Ophøjethed ved
den Indrømmelse, Platon gør ved at anerkende, at ogsaa
Smudset maa have sin Ide. Ideen om Smuds, siger Dr.
Ræder, kan dog ikke have været blandt de Ideer, som
Guder og Mennesker ifølge »Faidros« drog ud for at skue.
Hertil er at sige, at Ideen om Smuds ikke selv er smudsig.
Smudsets Begreb eller Ide vilde være den Samfoldighed af
Elementer og Love (kemiske og mekaniske Love, som vi nu
vilde kalde dem), der betinger, hvad der for os bliver
% --- p. 11 ---
\setcounter{footnote}{0}%
\opage{11}Fænomenet Smuds. Naar Platon selv har følt Betænkelighed
ved dette Begreb, er det, fordi det ikke er et rent
intellektuelt Begreb, men tillige et Vurderingsbegreb: Smudset
betegner noget, som man afskyer, vender sig fra eller
skaffer afvejen. Men det kan meget vel omfattes med rent
intellektuel Interesse, idet man kan se helt bort fra Værdisynspunktet.
Dette er ikke andet, end hvad Platon selv tidligere
har hævdet. I »Staten« (p. 476 A) erklæres udtrykkelig,
at de sande Filosofer ikke blot danner sig Begreber
om Retfærdighed, men ogsaa om Uretfærdighed, ikke blot
om det Gode, men ogsaa om det Onde. Altsaa har han
selv her anerkendt en Ide om det moralske Smuds. Det
er da ogsaa fra denne Side set en Inkonsekvens, naar han
en Stund er vegen tilbage fra at anerkende Smudsets Ide.
Han maatte konsekvent sige med Herakleitos: »Ogsaa her
er Guder!«

b. Naar Platon ikke har været klar over Forskellen
mellem de forskellige Arter af Ideer, kommer det af, at
han ikke har sondret mellem de forskellige Veje, ad hvilke
han kommer til de forskellige Ideer. Der er to Former af
Tankearbejde, Platon anvender uden bestemt at skelne
mellem dem.

Den ene bestaar i en Tilbagegaaen, en Regres fra foreliggende
Emner til deres Forudsætninger. Den fremtræder
tydelig i den Rolle, Begrebet Lighed, især i dets ideale
Form som Identitet, spiller hos ham. Parmenides havde
allerede set dette Begrebs fundamentale Betydning, Parmenides’
»Ene« betyder det, der er i fuldkommen Identitet med
sig selv, ensartet i alle Henseender og derfor kontinuerligt
(fr. 8. Diels.). Denne Erklæring har man træffende kaldt
en Akt, ved hvilken den menneskelige Fornuft tog sig selv
% --- p. 12 ---
\setcounter{footnote}{0}%
% CHECK p. 12: notes paired with the marks in order (the layer misread a number)
\opage{12}i Besiddelse, idet den hævder Tankens øverste Lov\footnote[1]{\textsc{Victor Brochard}: \textit{Études de Philosophie ancienne et de Philosophie moderne,} p. 133. --- Smlgn. ogsaa min Afhandling \textit{Parmenides fra Elea} (Mindre Arbejder. Anden Række) p. 141 f. (Modsætningen mellem, hvad vi nu kalder formel og real Videnskab fremtræder for første Gang i principiel Form hos Parmenides).}. Den
var paa én Gang en Protest mod Mangfoldigheden i Gudeverdenen
og mod det Kaos, menneskelig Erfaring syntes at
frembyde. Den sætter den skarpe Forskel mellem den indre
Tankesammenhæng, der tilsidst hviler paa Identitetsprincipet,
og Erfaringens mangfoldige Emner, der synes at
kunne forbindes paa saa mange forskellige Maader. Men
hos Parmenides skelnedes der endnu ikke mellem Enhed,
Identitet, Kontinuitet, Ensartethed og Uforanderlighed; alle
disse Udtryk stod som Betegnelser for det samme. Zenon
supplerede, som allerede berørt, sin Mesters Tankegang ved
sin Polemik mod Flerhed, Diskontinuitet og Bevægelse.
Fra mathematisk Side betragter man da ogsaa nu Zenon
som den, der »med sit skarpe Blik har peget hen paa de
Vanskeligheder, som Matematiken maatte overvinde for at
blive eksakt Videnskab«\footnote[2]{\textsc{Zeuthen}: \textit{Hvorledes Mathematiken i Tiden fra Platon til Euklid blev rationel Videnskab} (Vidensk. Selsk. Skr. 1917) p. 15. Se ogsaa \textit{Sur les livres arithmétiques d’Euclid} (Vid. Selsk. Overs. 1910), p. 432.}. Det var i Kraft af Identitets- og
Kontinuitetsidealet, som dogmatisk fastholdtes, at Zenon
med saa stor Energi betonede Umuligheden af en rationel
Opfattelse, naar man forudsatte Flerhed og Diskontinuitet.
Overfor denne dogmatiske Opfattelse har Platon den Fortjeneste
at have hævdet Nødvendigheden af flere Tanker
end den ene, hvis man vil forstaa Verden. Det er Betydningen
af, at han sætter Ideverdenen istedetfor det parmenideiske
Ene. Og en ganske særlig Fortjeneste har han ved
at undersøge de Veje, ad hvilke vi kommer til de Tanker,
der gør os Tilværelsen forstaaelig. Herhen hører den alle%
% --- p. 13 ---
\setcounter{footnote}{0}%
\opage{13}rede omtalte Tankegang i »Faidon«, hvor det vises, at
Maalestokken for Lighed ligger i Tilnærmelse til »Lighedens
Ide«, den rene Identitet. Herhen hører fremdeles
Tankegangen i Slutningen af »Statens« sjette Bog, hvor
han fra Iagttagelsernes Mangfoldighed og Foranderlighed,
der i det højeste gør Tro og Formodning mulige, stiger op
til matematiske Love og fra dem igen til de Forudsætninger,
paa hvilke de, og dermed al sand Videnskab, hviler, Forudsætninger,
som det er den egentlige Filosofis (Dialektikens)
Opgave at fremdrage for at give dem en afsluttende Begrundelse.
Gennem saadan analytisk eller regressiv Tænken
naar han til det Stade, paa hvilket Parmenides fra første
Færd havde stillet sig, og som Zenon havde søgt at forsvare
gennem Kritik af det modsatte Stade. Men Platonismen
og Eleatismen stemmer saa tilsidst overens i, at det
sidste Grundlag for at Tænkning er et Værende, --- ikke
blot Forudsætninger og Principer for vor Erkendelse, men
selve Tilværelsens eget Væsen. Ideverdenen er den højeste
Form for Eksistens, ligesom Parmenides’ »Ene« var den
eneste Form for Eksistens.

Den anden Form, under hvilken Platon’s Tankearbejde
fremtræder, er Ordning og Inddeling. Det er ad denne Vej,
han --- som man ser af »Parmenides«, efter nogen Betænkelighed
--- kommer til at opfatte Ideerne som Artsbegreber,
der hvert for sig gælder for en Gruppe af Emner.
Denne Vej gaar han f. Eks. i »Statens« fjerde Bog, naar
han skelner mellem forskellige »Dele« af Sjælen og bygger
denne Inddeling af sjælelige Livsytringer paa disses forskelligartede,
endog modsatte eller modstridende Karakter.
I »Faidros« stilles saa den bestemte Fordring, at man
baade skal kunne sammendrage, hvad der er fælles for de
forskellige Emner, og kunne følge, hvorledes dette Fælles
% --- p. 14 ---
\setcounter{footnote}{0}%
\opage{14}optræder under forskellige Former. Den herpaa byggede
Inddeling skal gennemføres, indtil man standser ved det
aldeles enkelte, hvor ingen Karakteristik længere er mulig.
I senere Dialoger (»Sofisten« og »Statsmanden«) beskriver
han gennem Eksempler fra det daglige Liv og i bred og
vidtløftig Udførelse den Inddelingens Kunst, der baade
lader Ligheder og Forskelligheder komme til deres Ret. De
Artsbegreber, som dannes paa denne Maade, var for Platon,
ligesom de logiske og etiske Begreber, ideale Eksistenser,
evige og uforanderlige Væsener. Og ligesom alle enkelte
Ligheder i Verden tilsidst fandt deres Forklaring ved
Lighedens evige Ide, saaledes indeholdt ogsaa Artsbegreberne
Forklaringen af de Emner, for hvilke de gælder.
Derved forstaas, at Platon ikke havde Grund til at skelne
skarpt mellem de forskellige Arter af Ideer. Hvad der for
ham var afgørende, var, at de alle indeholdt de sidste
Grunde til, at Tilværelsen faktisk er, som den er. Betegner
vi ved Ordet Ontologi den Tankegang, ifølge hvilken
de Tanker, vi former under vort Erkendelsesarbejde, umiddelbart
og uvilkaarlig staar som evige og forudsætningsløse
Væsener, saa er Platon’s Idelære Ontologi ligesaa vel som
Parmenides’ »Ene«.

c. Vi vender nu tilbage til Dialogen. --- Den første Indvending,
Parmenides retter mod Idelæren, efter at han har
faaet Platon til at opfatte denne i hele dens Konsekvens,
angaar Forholdet mellem Ideer og Emner. Dette Forhold
havde Platon beskrevet som en Delagtighed i Ideerne fra
Emnernes Side. Alle Ting, som ligner hverandre, er derved
delagtige i Lighedens Ide; alle Dyr delagtige i Dyrets Ide
o. s. v. Men, spørger Parmenides, hvorledes kan en Ide
bestaa som den Enhed, den er, naar en Mangfoldighed af
Emner skal have Del i den? Hvert Emne faar da kun en
% --- p. 15 ---
\setcounter{footnote}{0}%
\opage{15}Del af den hele Ide! Men Ideen kan ikke udstykkes, ligesaa
lidt som den helt kan være tilstede i noget enkelt
Emne (p. 131).

Platon indrømmer Indvendingens Rigtighed. Og at den
virkelig rammer, ligger deri, at ligesaa almen, evig og uforanderlig
som Ideen er, ligesaa konkrete, forgængelige og
skiftende er Emnerne. Denne Modsætning kan kun overvindes,
naar det indses, at Ideen, Almenbegrebet, maa
fremtræde paa en særlig, individuel Maade i hvert enkelt
Emne. Der er da ikke Tale om udvortes Deling eller om
Fællesskab under samme Dække. Baade Retfærdighed, Lighed,
Menneske, Ild og Smuds fremtræder i vor Erfaring
stedse i individualiserede Former. Det var allerede Aristo
teles opmærksom paa, og det blev ivrig drøftet i Middelalderens
Skolastik og i den nyere Filosofi (Spinoza, Leibniz
og Berkeley). Det er Individualitetsproblemet, der gør
sig gældende. Platon’s Stordaad var det at hævde de almene
Tankebestemmelsers Betydning for vor Erkendelse.
Det næste Skridt, som først kunde gøres efter lange Tiders
Drøftelser, var at stille saadanne Tankebestemmelser i det
rette Forhold til de Emner, for hvilke de skal passe. Pla
ton stiger med Begejstring op til sine Ideer, men bryder
sig saa derefter ikke om, hvorledes det gaar med de Emner,
han har brugt som Udgangspunkter og Trin ved sin
Opstigen. Intetsteds, hverken i selve Dialogen »Parmenides«
eller i andre Dialoger besvarer Platon direkte den Indvending,
som her er kommen til Orde, og som paa et vist
Trin af hans Tankeliv maa være fremtraadt for ham. Men
hans Tanker kredser stadig om den, og det vil vise sig,
at det gaar op for ham, at logisk og mathematisk Tænkning
ikke blot kan aabne os Adgang til de rene Ideers
Verden, men ogsaa kan gøre det muligt at finde Tanke%
% --- p. 16 ---
\setcounter{footnote}{0}%
\opage{16}udtryk for Erfaringsemnernes skiftende Former, deres Tiltagen
og Aftagen, deres Svingninger og deres Forandringer.
Ved den tredie Indvending vil vi komme tilbage til dette
Punkt.

Parmenides’ anden Indvending gaar ud paa, at naar
Ideer skal være nødvendige til Forstaaelse af Emnerne,
fordi disse ellers fremtræder som en usammenhængende
Mangfoldighed, saa behøves der nye Ideer til Forstaaelse
af Forholdet mellem hine Ideer og de Emnegrupper, de
svarer til. Istedetfor Emnernes usammenhængende Mangfoldighed
møder nu den usammenhængende Mangfoldighed
af Ideer og Emner. Forskellen mellem Ide og Emne maa
stille det samme Problem som Forskellen mellem Emne
og Emne. Vi maa f. Eks. søge en Ide, der forklarer Forholdet
mellem Størrelsens Ide og de store Ting. Der maa
jo være noget, som er fælles for Størrelsen og de store
Ting. Det bliver en Ide i anden Potens. Men dermed er
vi ikke færdige, thi Spørgsmaalet bliver, hvorledes Forholdet
er mellem Ideen i anden Potens paa den ene Side,
Ideen i første Potens og Emnerne paa den anden Side. Vi
kommer altsaa ind i en uendelig Række og bliver aldrig
færdige (p. 132 A, B).

Denne Indvending omtales ofte hos Aristoteles og har
ført til et teknisk Udtryk: »det tredie Menneske«. Aristoteles
forudsætter bekendt, hvad der menes dermed. Ligesom
Mennesket, som vi kender det i Erfaringen, forudsætter
Menneskets Ide (det Begreb om, hvad et Menneske er, der
gælder for de enkelte Menneskeeksemplarer), saaledes maa
vi i Kraft af Idelæren antage en ny Menneskeide, der gælder
for Forholdet mellem den første Menneskeide og
Menneskeeksemplarerne --- altsaa et tredie Menneske, og
dette kan fortsættes i det uendelige. Denne Indvending skal
% --- p. 17 ---
\setcounter{footnote}{0}%
% PRINTED AS IS: „Aristotetes’“ (for Aristoteles’), p. 17, seen at the image (blind read b1-02-03)
\opage{17}(ifølge et Scholie af Alexandros fra Afrodisias til Aristotetes’
Metafysik) være fremsat af Sofisten Polyxenos og har
sandsynligvis været Platon bekendt, da han skrev »Parmenides«\footnote[1]{Se \textsc{Baeumker}: \textit{Über den Sophisten Polyxenos.} (Rheinisches Museum. N. F. XXXIV).}.
Maaske har han allerede i »Staten« X p. 597 C
forsøgt at besvare den. Han siger her, at Gud kun har
skabt en eneste Ide for hver Emnegruppe, »hvad enten
han ikke vilde frembringe flere, eller der forelaa en Nødvendighed
for kun at frembringe en eneste«. Platon vil
have en Ende paa det, lukke for en Tankegang, der vilde
kunne fortsætte sig i det uendelige; men det lykkes ham
kun ved et teologisk Postulat. Ligesom senere for Aristoteles
(og endnu senere for Hegel) er den Konsekvens, at
en Opfattelse fører ud i en uendelig Tankerække, for ham
en tilstrækkelig Gendrivelse af en saadan Opfattelse. Det er
i Virkeligheden den samme Tankegang, der ligger til Grund
ved Zenon’s Kritik af Flerhed og Bevægelse.

For den nyere Tids Erkendelsesteori staar det derimod
som en naturlig Ting, at Erkendelsesarbejdet ikke kan afsluttes.
Selv de sidste Forudsætninger, til hvilke Analysen
af vor Erfaring fører os tilbage, tør vi ikke betragte som
absolute og »i sig selv grundede«. Der er Tankerækker,
der stadig kan fortsættes, og indenfor hvilke der stadig vil
kunne findes flere Mellemled, saa stor Sammenhæng der
end hidtil er funden. Muligheden af en Rækkes stadige
Fortsættelse, hvormange Led og Mellemled der end findes,
indeholder for os ingen Modsigelse og ingen Vanskelighed.
Vi holder os til den Lov, ifølge hvilken den stadige Fortsættelse
er mulig, og i denne Lov ser vi Erkendelsens Maal.
Og naar det drejer sig om Forudsætninger, der er nødvendige
for at vinde Forstaaelse, ser vi i Muligheden af stadig
% --- p. 18 ---
\setcounter{footnote}{0}%
\opage{18}Fortsættelse i Kraft af samme Lov et Bevis for, at vi her
staar ved en af den menneskelige Tænkeevnes Ejendommeligheder.
De logiske Grundsætninger, som forudsættes ved
ethvert Forsøg paa at begrunde en Antagelse, kan ikke selv
begrundes, da ethvert Forsøg derpaa vilde forudsætte dem.
Paa tilsvarende Maade har man i den Forudsætning, at
hvad der gælder for Størrelsen n ogsaa gælder for n + 1,
fundet den fundamentale Ejendommelighed ved matematisk
Tænken\footnote[1]{Smlgn. herom \textit{Totalitet som Kategori.} (Vid. Selsk. Skr. 1917), p. 34---36.}, særlig i dens Forskel fra formel logisk Tænken.
--- Allerede fra den nyere Tids Begyndelse var store Filosofer
klare over, at vi i Rækker, der stedse kan fortsættes
efter samme Lov, maa se Udtryk for et Grundforhold for
vor Erkendelse, maaske for selve Tilværelsen. Dette fremtræder
tydelig hos Bruno og Spinoza\footnote[2]{Se \textit{Spinoza’s Ethica.} (Vid. Selsk. Skr. 1918), p. 8---9; 38 f.}.

\textsc{Wilamowitz-Moellendorff}, som i det hele ikke tillægger
Indvendingerne mod Idelæren stor Betydning, mener,
at den anden Indvending, den jeg her drøfter, let kan afvises
ved følgende Betragtninger. Hvad der er fælles for
flere Ting, findes ved at se bort fra Forskellighederne mellem
disse Ting og holde sig til Lighedspunkterne. Ideen
dannes altsaa ved Subtraktion. Men naar man saa igen
sammenligner Ideen med Tingene, er der ikke mere noget
at trække fra, altsaa ingen Mulighed for en ny Ide\footnote[3]{\textit{Platon} II, p. 26.}. Hertil
maa svares, at netop naar der ikke er mere at trække fra,
staar vi ved et saadant Forskellighedernes Kaos, som det
var Idelærens Opgave at hæve os over. En Opgave vilde
der altsaa i hvert Tilfælde blive tilbage. Men Hovedsagen
er, at det ikke er i Platon’s Aand at betragte den Tankeproces,
der fører til Ideer, som en mekanisk isolerende
% --- p. 19 ---
\setcounter{footnote}{0}%
\opage{19}Fremgangsmaade. Idedannelsen er Frugt af en sammenlignende
Undersøgelse, der fører til en Syntese, eller som
Platon kalder det en Synopsis, en Samskuen, der fører de
mangfoldig adspredte Fænomener tilbage til én Tanke.
(Faidros p. 265 C: εἰς μίαν ἰδέαν συνορῶντα ἄγειν τὰ πολλαχῇ
διεσπαρμένα). Og man vil ikke kunne komme fra, at en
ny Sammenligning kræves overfor Indholdet mellem den
saaledes dannede Ide og de derfra forskellige Emner, ligesom
der krævedes en Sammenligning overfor Forholdet
mellem de indbyrdes forskellige Emner.

\textsc{Jackson}\footnote[1]{Journal of Philology X p. 298; 334.} er klar over, at Platon’s Idelære, saaledes
som han har fremstillet den i »Faidon« og i »Staten«,
rammes af den her drøftede Indvending. Men han lægger
% PRINTED AS IS: „for-|forsøger“ (the syllable set twice), p. 19, seen at the image
megen Vægt paa en Udvej, som Platon selv et Øjeblik forforsøger,
den nemlig, at man opfatter Forholdet mellem
Ideer og Emner saaledes, at hine er de evige Typer, Forbilleder
for disse, saa at det, der blev kaldt Delagtighed i
Ideen, i Virkeligheden er Lighed med Ideen (p. 132 D).
Derved slipper man, mener Jackson, for den mekaniske
Opfattelse, ifølge hvilken Ideen ligesom skulde udstykkes
blandt Emnerne. Jackson tager nærmest Platon’s teologiske
Løsning fra »Statens« 10. Bog op. Naar der ved Frembringelsen
af Emnerne stod visse evige og uforanderlige
Typer som ledende Forbilleder, saa vil det gøre Ordning
og Inddeling af Emnerne mulig, og Studiet af disse vil
føre os nærmere til Erkendelsen af Ideerne som de fuldkomne
Udtryk for, hvad Emnerne kun i Tilnærmelse frembyder.


Det er muligt, at denne Opfattelse af Ideerne som Forbilleder
i Grunden var den, Platon i sit Inderste hyldede.
I den kunde, synes det, hans filosofiske og hans digteriske
% --- p. 20 ---
\setcounter{footnote}{0}%
\opage{20}Trang saavel som hans religiøse Trang finde Tilfredsstillelse
paa én Gang. Men Parmenides gør (p. 132 E) med
fuld Ret opmærksom paa, at Vanskeligheden kommer
igen ganske paa samme Maade som før, da Lighedsforholdet
mellem Forbillede og Emne ligesaa vel maa føre
til Antagelse af en Ide, som Lighedsforholdet mellem Emnerne
indbyrdes. Ogsaa her staar vi overfor den uendelige
Række, som den græske Tænkning nærede saa stor
Afsky for. ---

En anden Vej har Platon et Øjeblik grebet. Der var
den Mulighed, at Ideerne kun er vore Tanker, kun boer
i Sjælene, intet andet Sted. Saa vilde man, synes det, slippe
for alle betænkelige Konsekvenser med Hensyn til Ideernes
Indhold til Emnerne. Men denne Mulighed afskærer Parmenides
ligestraks, og med fuld og hurtig Tilslutning fra
Platon’s Side, idet han hævder, at enhver Tanke jo dog
maa være en Tanke om noget, og det et noget, som har
Væren. En Tanke, der ikke er en Tanke om noget værende
(νόημα ὄντος), er en Meningsløshed (132 B, C). --- Det er
en i høj Grad karakteristisk Udtalelse, der kaster Lys over
Idelæren og over antik Tænkning idethele. Man ser her,
hvor nær Ontologien laa for den begyndende Tænkning.
Tankerne kan komme og gaa, men Tankeindholdet maa
have evig og objektiv Bestaaen.

Nu har man ganske vist i den nyeste Tid villet nægte,
at Platon’s Idelære har denne ontologiske Karakter. Man
har paastaaet, at Ordet »Væren« hos Platon, naar det bruges
om Ideerne, kun betyder Indhold af et Udsagn\footnote[1]{»Nie und nirgends bedeutet »Sein« bei Platon, wenn von der Idee ausgesagt, etwas anderes als den Aussageinhalt«. \textsc{Natorp}: \textit{Platon.} (Grosse Denker. Herausg. von Aster. I.) p. 111. --- Allerede i sit Skrift \textit{Platon’s Ideenlehre} (1903) har Natorp udtalt denne Opfattelse, til hvilken J. A. \textsc{Steward} (\textit{Platon’s Doctrine of Ideas} p. 77) har sluttet sig.}. Og
% --- p. 21 ---
\setcounter{footnote}{0}%
\opage{21}man har paa Grundlag af denne Paastand ment, at medens
den Indvending mod Idelæren, der kaldtes »det tredie
Menneske«, var uafviselig, naar Ideerne opfattedes som
værende Ting, der stod i Modsætning til Emnerne som
andre værende Ting, falder denne Indvending ganske bort,
naar man indser, at Platon’s Ideer kun er metodologiske
Synspunkter, der leder til gyldig Erkendelse.

Det er den moderne Erkendelsesteoris, af Kant grundlagte
Synsmaade, man ved denne Opfattelse vil lægge ind
hos Platon. I en vis Forstand er det heller ikke urigtigt,
at Ideerne for Platon betyder metodologiske Synspunkter.
Det er ud fra Ideerne, han betragter alle Emner og fælder
sin Dom om dem, baade hvad deres Realitet og hvad deres
Værdi angaar. Det er --- som allerede ovenfor antydet ---
hans Betydning i Tænkningens Historie at have indset, at
al vor Erkendelse beror paa Forudsætninger, det gælder
om at blive sig bevidst. Han vilde i sin Idelære paa sin
Maade fyldestgøre sin Mesters Fordring om Selverkendelse.
Men Sagen er netop, at det, at Ideerne var metodologiske
Synspunkter, for ham ingenlunde udelukkede, at de ogsaa
var værende Væsener. Der var her for ham intet Enten-Eller,
men netop et nødvendigt Baade-Og. Han skelnede
ikke mellem Gyldighed og Væren. Derfor kunde erkendelsesteoretisk
Interesse og ontologisk Dogmatisme gaa i Et hos
ham. Han staar her paa lignende Standpunkt overfor sine
Ideer, som Pytagoræerne overfor deres Tal og Eleaterne
overfor deres »Ene«. Ogsaa moderne Matematikere tillægger
undertiden de matematiske Grundtanker ontologisk Betydning,
uden at det hindrer dem i Anvendelse af exakte
matematiske Metoder.

For Kant var der her et Enten-Eller. Det førte ham til
% --- p. 22 ---
\setcounter{footnote}{0}%
\opage{22}den Erklæring, at det stolte Navn »Ontologi« maa afløses
af det beskednere »Forstandens Analytik«.

Men selvom det var berettiget at lægge den kritiske
Filosofis Synsmaader ind hos Platon, vilde »det tredie
Menneske« ikke forsvinde. Thi der maa opkastes det Spørgsmaal,
hvilket Forholdet egentlig er mellem vore metodologiske
Synspunkter og de Emner, de anvendes paa. Finder
Emnerne, med alle deres kvalitative og temporale Forskelligheder,
deres fuldkomment udtømmende Udtryk i de
Love og Ligninger, som Anvendelsen af vore metodologiske
Synsmaader kan føre til? Er der her et Forhold af
Identitet eller blot af Analogi? Bestaar der i hvert Tilfælde
ikke stedse et Problem her, af samme Art som det, der
maatte vedblive at bestaa for Platon i Forholdet mellem
Ideer og Emner? Naar Natorp nægter dette sidste Problems
Tilstedeværen, er det, fordi han, som Repræsentant for
Marburgerskolen\footnote[1]{Smlgn. min Kritik af Marburgerskolen i \textit{Totalitet som Kategori} p. 46---48. (Vidensk. Selsk. Skr. 1917).}, paa det afgørende Punkt forveksler
Identitet og Analogi, ligesom allerede Platon selv har
gjort det.

For enhver Verdensanskuelse --- teologisk, filosofisk eller
naturvidenskabelig --- rejser der sig her et Problem, som
stedse stiller sig paany, og som opstaar ved, at der nu
engang er Forskel paa Iagttagelse (Oplevelse) og Tydning.
Det kan være, at Iagttagelse med Nødvendighed fører til en
vis Tydning; men saalænge man ikke kan vende Forholdet
om og ud fra de Tanker, Tydningen indeholder, udlede
netop de bestemte Iagttagelser, man byggede paa, saalænge
kan det ikke bevises, at et Identitetsforhold er tilstede.
Det er umuligt at udlede Mangfoldighed af Enhed, medens
Analyse af Mangfoldigheden meget vel kan begrunde An%
% --- p. 23 ---
\setcounter{footnote}{0}%
\opage{23}tagelsen af en Enhed (i Betydning af lovmæssig Sammenhæng).
I teologisk Tænken optræder Problemet Enhed—
Mangfoldighed i Forholdet mellem Gud og Verden; men vi
faar derved kun samme Problem igen --- blot i større
Dimensioner\footnote[1]{Smlgn. min \textit{Religionsfilosofi} § 17, hvor Parmenides’ Indvending allerede er anerkendt og paavist i dens Betydning for Teologien.}. ---

Vender vi tilbage til Parmenides, saa har han sin største
Indvending tilbage, --- --- som dog vil vise sig at hænge
nøje sammen med de to første, at være som en Fortsættelse
af dem. --- Naar Ideerne er som en Verden for sig,
med deres egne Love (αὐταὶ πρὸς αὐτὰς τὴν οὐσίαν ἔχουσι),
saa vil de ikke kunne erkendes af Mennesker, der høre
hjemme i den sanselige Erfarings Verden, hvis Ting ligeledes
bestemmes ved deres egne indbyrdes Forhold (αὐτὰ
πρὸς αὐτά ἐστιν). Den, der er Herre, staar f. Eks. i Forhold
til sine Slaver, ikke til Begrebet (Ideen) Slave, og omvendt,
hvorimod Begrebet (Ideen) Herre og Begrebet (Ideen)
Slave staar i Forhold til hinanden. Og saaledes forholder
Begrebet (Ideen) Videnskab sig til Begrebet Sandhed (Sandhedens
Ide)\footnote[2]{Smlgn. hermed \textit{Faidros} p. 247 E (I Præexistensen skuer Sjælene absolut Videnskab: --- τὴν ἐν τῷ ὅ ἐστιν ὄν ὄντως ἐπιστήμην οὖσαν).}, men vor Videnskab forholder sig til den
Sandhed, vi kan naa. Ideerne bor jo ikke i os, men i deres
egen Verden, og Ideernes Verden kan ikke gøre os en fra
den forskellig Verden forstaaelig. End ikke Gud, som dog
maa forudsættes at staa i Forhold til Videnskabens Ide\footnote[3]{At Gud er Gud, skyldes hans Forhold til Ideerne. Smlgn. \textit{Faidros} p. 249 C (Hvad Sjælene saa i Præexistensen, Ideerne, er det, πρὸς οἷσπερ θεὸς ὤν θεῖος ἐστιν).},
vilde kunne erkende den Verden, vi lever i, og ligesaalidt
vil vi fra vor Verden kunne stige op til Ideernes Verden.
% --- p. 24 ---
\setcounter{footnote}{0}%
\opage{24}Hverken Afbilledlighed eller Delagtighed vil der her kunne
være Tale om (p. 133 B—134 B).

At denne Indvending rammer, kan ses af den Erkendelsesteori,
Platon selv har udviklet i »Statens« sjette Bog.
Indenfor Tankens Verden (τὸ νοητόν) skelnes mellem
Matematiken, der arbejder med visse Forudsætninger, hvis
Grunde ikke angives, og Dialektiken, der kun benytter
Hypoteser som Trin og Tilløb til at naa det Forudsætningsløse,
Principet for alt, og denne højeste Erkendelse skrider
frem fra Ide til Ide og ender i Ideer (Staten VI p. 511 B).
Idelæren er altsaa et sluttet Hele, der ikke peger ud over
sig selv. Har man blot fundet Indgangen til den, gaar man
frem ad den rene Deduktions Vej, holdende sig til hvad
der følger af Principet (ἐχόμενος τῶν ἐκείνης [τῆς ἀρχῆς]
ἐχομένων), og man ender i Ideer (τελευτᾲ εἰς εἴδη) og
kommer ikke ud af Ideernes Verden. --- Det staar her som
uforstaaeligt, hvorledes denne rene Deduktion kan faa nogen
Betydning for Erfaringens Verden, hvad dog Platon udtrykkelig
forudsætter, da han lægger Vægt paa, at kun de,
der har opnaaet denne højeste Erkendelse, egner sig til at
lede de menneskelige Anliggender. Men han har ikke vist
nogen Mulighed af, at den rene Deduktion kunde komme
i nogen Vexelvirkning med den Kundskab, der vindes ved
Iagttagelse.

Parmenides’ tredie Indvending drager nu netop kun
Konsekvensen af Idelæren som rent deduktiv Erkendelse.
Men han gaar et Skridt videre, idet han nægter Muligheden
af at naa det Punkt, hvor denne Deduktion kunde begynde.
End ikke de af Platon nævnte Trin og Tilløb bliver mulige,
idet Erfaringens Verden ogsaa er afsluttet i sig selv.
Hvor skulde det Punkt ligge indenfor Erfaringens Verden,
der kunde gøre en Opstigen, et Spring op i Ideverdenen
% --- p. 25 ---
\setcounter{footnote}{0}%
\opage{25}mulig? --- Tankegangen er den samme, som den Kierkegaard
anvendte overfor Hegel\footnote[1]{Smlgn. »\textit{Søren Kierkegaard som Filosof}«\textsuperscript{2} p. 61 f.}, idet han fordrede, at »Begyndelsens
Dialektik« skulde gøres klar, før man begyndte
paa Systemet. Hvorledes gøres Overgangen fra Analyse af
Erfaring til den deduktive Syntese? Kierkegaard drager heraf
den Slutning, at der ingen forudsætningsløs Filosofi gives. ---

Allerede de to første Indvendinger berørte væsentlige
Sider ved Idelæren. Men den tredie Indvending drager
Grunden helt bort under den. Det tjener til Platon’s Ære,
at han har ladet den radikale Betænkelighed ved Idelæren
komme saa klart til Orde. Man kan sige, at han her er
kommen lige til det Punkt, hvor Kant i Aaret 1772 fandt
Grundlaget for sin afgørende Problemstilling. I sit Brev til
Marcus Hertz 21. Febr. 1772 spørger Kant netop, hvorledes
det kan være, at de Begreber, vi danner ved vor aandelige
Virksomhed, kan faa Betydning for en Verden, vi ikke
selv har frembragt\footnote[2]{Se \textit{Kontinuiteten i Kant’s filosofiske Udviklingsgang} (Vid. Selsk. Skr. 1893) p. 52.}. Hos Platon formes Spørgsmaalet saaledes:
hvorledes kan menneskelig Videnskab blive guddommelig
Videnskab? ---

Ligesaa lidt som de første Indvendinger fører denne
tredie Indvending Platon til at opgive Idelæren, og vi vil
straks komme til at se, hvorfor han ikke opgiver den.
Men denne tredie Indvending fører ham dog til at foretage
Ændringer i sin Filosofi. Det er Forholdet mellem Deduktion
og Induktion, der er det afgørende. Disse to Metoder
gaar hos Platon i de store Dialoger (Symposion, Faidon
og Staten) hver sine Veje. Det maatte gælde om at finde
Mulighed for en Samvirken mellem dem. Paa en saadan
Samvirken bygger den moderne Videnskab fra Renais%
% --- p. 26 ---
\setcounter{footnote}{0}%
\opage{26}sancen af, og som vi vil komme til at se, bevæger Platon’s
senere Tænkning sig ogsaa i denne Retning, selv om
den ikke naar Maalet. ---

Efterat de tre Indvendinger er drøftede, siger Parmenides,
at saadanne Vanskeligheder kommer man til, naar
man tillægger Ideerne Væren. Men hvis man paa Grund
af disse Indvendinger vil opgive Idelæren, saa vil man
ikke kunne give sin Tænken nogen bestemt Retning
(οὐδὲ ὅποι τρέπει τὴν διάνοιαν ἕξει) og Evnen til tankemæssig
Drøftelse (ἥ τοῦ διαλέγεσθαι δύναμις) vil saa til
sidst forsvinde (p. 135 ABC). --- Stærkere og mere karakteristisk
kunde det ikke udtrykkes, at Ontologien maatte
fastholdes, og hvorfor den maatte fastholdes. Parmenides
og Platon er her fuldt enige. Og Parmenides er overhovedet
i denne Dialog Platon selv, eller en Del af ham. Man
kan tale om den platoniske Parmenides ligesom om den
platoniske Sokrates.

\chaphead{2. Dialogens anden Del.}


Da den platoniske Parmenides’ unge Ven (den platoniske
Sokrates) er i Forlegenhed med, hvad der nu skal
blive af Filosofien, anerkender den gamle Filosof først
hans Interesse for Tankearbejde, i Modsætning til dem, der
lader Tanken strejfe om uden Lov og Regel. Og derpaa
opfordrer han ham til at øve sin Tanke ved at gaa ud fra
skarpt bestemte Forudsætninger. For at finde, hvor Problemerne
virkelig ligger, maa man baade undersøge, hvad
der følger af en vis bestemt Forudsætning, og hvad der
følger af den modsatte. Nu staar Enhed og Flerhed, Lighed
og Ulighed, Hvile og Bevægelse, Tilblivelse og Opløsning,
Væren og Ikke-Væren overfor hinanden som Modsætninger.
Vi maa undersøge, hvad Følgen bliver af at holde
% --- p. 27 ---
\setcounter{footnote}{0}%
\opage{27}sig til et enkelt Led indenfor ethvert af disse Modsætningspar.
Kun ved en saadan alsidig Drøftelse kan man opnaa
Sandhed og erhverve Fornuft (135 C—136 E). Parmenides
vil da nu i det følgende drøfte sin egen Forudsætning om
det Ene som Forudsætning for alt (137 B).

Den af Parmenides foreslaaede Problembehandling har
tydelig Foranledning i de foregaaende Undersøgelser. Idelærens
Vanskeligheder viste sig overalt at bero paa Modsætningen
mellem Ide og Emner. Eleaterne slap for dette
Problem, idet de paa Forhaand slog fast, at kun »det Ene«
(svarende til Platon’s Ideverden) har Realitet. Det er en
stor Indrømmelse, Platon lader Parmenides gøre, naar han
lader ham gaa ind paa at drøfte baade den Forudsætning,
at »det Ene« er, og den, at det ikke er. Dog var denne
anden Forudsætning, den, at »det Ene« ikke er, at altsaa
det Mangfoldige er, allerede bleven drøftet af Zenon i hans
berømte Skrift, hvis Tankegang da ogsaa spores paa mange
Steder i det følgende. Men Platon lader Parmenides indføre
nye Distinktioner, som strider mod de eleatiske Forudsætninger
og gør Drøftelsen meget indviklet, for ikke at
sige helt dunkel. Allerførst maa vi sige, at Emnet for den
foreslaaede Metode ikke var heldig valgt. Det var jo Platonismen,
ikke Eleatismen, der skulde undersøges. Vel var
Platon enig med Parmenides i, at vore Tankers gyldige
Indhold maa betyde noget Værende, og ogsaa i, at Tanken
søger at forbinde sit mangfoldige Indhold saa inderlig som
muligt, altsaa finde Enhed og Sammenhæng i det. Men
Platon’s Fremskridt fremfor Parmenides bestod netop i, at
Platon med stor Energi og Begejstring baade fordrede og
beskrev det aandelige Arbejde, der kunde føre til at finde
denne Enhed, som Parmenides uden videre forudsatte. I
»Menon«, »Symposion«, »Faidon«, »Staten« og »Faidros«
% --- p. 28 ---
\setcounter{footnote}{0}%
\opage{28}havde han skildret Udviklingen fra den sanselige Opgaaen
i Emnernes Strøm til Tankelivets Enhed og indre Sammenhæng.
Og det havde vist sig, at »Enhed« betyder lovmæssig
Sammenhæng, Rationalitet eller Identitet. I »Parmenides«
p. 139 D siges udtrykkelig, at Enhed og Identitet er
det samme (ἥπερ τοῦ ἑνὸς φύσις, αὕτη δήπου καὶ τοῦ ταὐτοῦ),
og hermed stemmer Udtalelser i »Staten« (p. 533 C; 537 C)
og i »Faidros« (p. 265 DE; 266 B). Det vilde have været
baade naturligere og frugtbarere, om det havde været
Identitetsbegrebet, og som Modsætning dertil Forskelsbegrebet,
der var blevet Genstand for Drøftelse, for at
Konsekvensen af kun at lægge det ene eller det andet af
disse to Modsætningsled til Grund kunde blive fremdragen.
Problemet Enhed-Mangfoldighed er afledet i Forhold til
Problemet Identitet og Forskel. Hvad Zenon vilde tillivs
i sine berømte Aporier var, som allerede omtalt, egentlig
ikke Bevægelsen, men Flerheden, den absolute Mangfoldighed
(hvad jeg har kaldt den kaotiske Forskelsrække), som
udelukker Forstaaelse. Det er --- som Platon klart har set ---
Forstaaelsens Mulighed, der er Eleaternes egentlige Problem,
som Parmenides drøfter positivt, Zenon negativt. For Platon
vilde det have været naturligst at lade Drøftelsen dreje
sig om selve Forstaaelsens, altsaa Videnskabens Mulighed.
Derved at Enhedsbegrebet skydes frem foran Identitetsbegrebet,
bliver det mere Eleatisme end Platonisme, der
drøftes i det følgende, og derved bliver Drøftelsen baade
ufrugtbar og uklar. Det er i Grunden hverken Eleatisme
eller Platonisme, overhovedet intet klart Stade, der undersøges.


Efter ofte gentagen Syslen med Dialogens anden Del
har jeg for mit Vedkommende faaet et afgjort Indtryk af
Mangel paa naturligt Anlæg, og tillige af mangelfuld Gennem%
% --- p. 29 ---
\setcounter{footnote}{0}%
\opage{29}førelse i det enkelte og af Mangel paa naturlig Afslutning.
Det var jo dog bebudet, at en nøjagtig Gennemtænken af
de forskellige mulige Hypoteser skulde hjælpe os paa Gled
igen i vor Filosofi, naar det Haab, vi havde næret til Idelæren
i dens hidtilværende Form maatte opgives. Men der
gives ingensomhelst Antydning i den Henseende.

Som allerede omtalt ligger her en af Grundene til, at
mange en Tid var tilbøjelige til at erklære Dialogen for
uægte. De Uklarheder og Fejlslutninger, anden Del frembød,
vilde man ikke skrive paa Platons Regning. Andre Forskere
har opfattet anden Del som et i polemiske Øjemed
forfattet, af Sofismer vrimlende Ungdomsskrift; saaledes
Otto Apelt (Beiträge zur Geschichte der griechischen Philosophie
p. 59; 66). \textsc{Wilamowitz} slutter sig i sin Platonbiografi
for en Del til Apelt. Han betragter Dialogen som ufuldendt,
frembragt i en Periode af Platon’s Liv, hvor han efter
praktiske og teoretiske Skuffelser forsøgte nye Veje. For mit
Vedkommende anser jeg det ogsaa for sandsynligt, at den
sidste Haand ikke er lagt paa Dialogens anden Del, der
stikker saa stærkt af mod den omhyggelig gennemarbejdede
og tankerige første Del. Jeg formoder da, at Platon ved
nærmere Eftertanke er ført til at stille det Hovedproblem,
der nu fremstillede sig for ham efter den strenge Selvkritik,
paa en Maade, der baade var klarere og bedre svarer
til hans Forudsætninger. De følgende Dialoger har, som
jeg i et senere Afsnit skal vise, afgjorte Antydninger i Retning
af en saadan ændret Problemstilling.

Med Hensyn til Gangen i Dialogens anden Del vil jeg
henvise til de Forskere, der mener at have fundet et positivt
Indhold i den, som \textsc{Eduard} \textsc{Zeller} (Platonische Studien
p. 174—179), \textsc{Jackson} (Journal of Philology XI p. 306—
331) og \textsc{Ræder} (Platon’s philosophische Entwickelung
% --- p. 30 ---
\setcounter{footnote}{0}%
\opage{30}p. 309—316). Disse Forfattere har som almindeligt Resultat
af denne anden Del af Dialogen fundet den Erkendelse,
at Enhed og Mangfoldighed stadig forudsætter hinanden,
saa at man ikke udelukkende kan lægge et af de to Modsætningsled
til Grund for sin Tænken. Det er kun et af de
mange Eksempler paa, at vore Begreber staa i Korrelationsforhold
og derfor tjener til gensidig Belysning og Bestemmelse
af hverandre. Men denne Sandhed vilde Platon sikkert,
dersom han havde gennemarbejdet sit Udkast, kunne
have udviklet paa en klarere og naturligere Maade og i
nøjere Sammenhæng med sine Grundtanker.

Kun nogle enkelte Punkter skal jeg her fremdrage til
Karakteristik af Dialogens anden Del. Det er dels en skarpt
gennemført Tanke, dels en forunderlig Sofisme, dels et
betydningsfuldt Tankeglimt.

Parmenides begynder med at undersøge sin egen Grundtanke,
Tanken om »det Ene«. Hvis Begrebet »det Ene«
forstaas ubetinget, kan der slet ingen Prædikater tillægges
det. Det kan ikke kaldes helt, da det ingen Dele kan have,
og naar det intet Hele er, kan det hverken have Begyndelse
eller Ende, altsaa ikke være i Tiden og ikke undergaa
nogen Forandring. Det kan ikke have nogen Beliggenhed,
da der saa forudsættes noget uden for det. Det kan
ikke være i sig selv, da det saa vilde være baade større og
mindre end sig selv, og heller ikke være i noget andet, da
der saa vilde være noget andet foruden det. Det kan ligesaa
lidt være i Hvile (»staa«) som bevæges, ligesaa lidt
være identisk med sig selv som forskelligt fra sig selv.
Men det kan heller ikke være identisk med noget andet
eller forskelligt fra noget andet. (I denne Sammenhæng er
det, at det røbes, at Enhed og Identitet er det samme). Og
hvis Prædikatet Væren betyder Anvendelighed af et af de
% --- p. 31 ---
\setcounter{footnote}{0}%
\opage{31}nævnte Begreber, særlig af Tidsbestemtheden, saa kan man
ikke kalde det værende. Intet Prædikat kan altsaa anvendes
paa det, og derfor kan det hverken benævnes, sanses eller
tænkes (137 C—142 A).

Platon lader her Parmenides drage Grunden bort under
sin egen Lære. I sit Læredigt havde den historiske Parmenides
--- som allerede ovenfor anført --- tillagt sit »Ene«
flere Prædikater: det er Et med Væren; det staar i Ligheds-
og Kontinuitetsforhold til sig selv; det har Kredsform.
Det paradoxale er, at Begrebet Enhed (Identitet) netop er
opstillet af Eleaterne som Princip for Tilværelsens Forstaaelighed.
Og nu viser det sig, at det selv er uforstaaeligt!
--- Naar dette var paavist, kunde der have været Grund
til at undersøge selve Begrebet Forstaaelighed. Det er den
Vej, ad hvilken Platon har fjernet sig fra Eleatismen. Og
den førte ham til en Flerhed af Ideer, svarende til forskellige
Grupper af Emner. I Dialogens første Del var det jo
udtrykkelig hævdet, at der ingen Modsigelse er i, at samme
Emne kan henhøre under forskellige Ideer, altsaa have
forskellige Prædikater. Denne Tankegang kunde Platon
have ført videre (og han gør det i senere Dialoger); men
der hentydes slet ikke til den i Dialogens anden Del.

Med Hensyn til Ordningen indenfor Dialogen maa man
undre sig over, at Kritiken af »det Ene« ikke har fundet
sin Plads i første Del, som Led i første Indvending mod
Idelæren. Kritiken af det »Ene« faar nemlig ikke blot Betydning
overfor Eleatismen, men ogsaa overfor Platonismen,
da hver enkelt Ide staar som Enhed i Modsætning til den
Mangfoldighed af Emner, der hører ind under den. Tillige
har jo Platon (som »Staten« p. 507 og det syvende Brev
viser) aldrig opgivet Haabet om at kunne føre alle Ideer
% --- p. 32 ---
\setcounter{footnote}{0}%
\opage{32}tilbage til en eneste Grundidé, og her kommer da Problemet
Enhed-Mangfoldighed frem igen.

Men ikke blot overfor Eleater og Platonikere har denne
mærkelige Kritik af Begrebet »det Ene« sin Betydning.
Atter og atter fremtræder baade i Religionens og i Tænkningens
Historie Forsøg paa at samle alt i én Tanke eller
paa at gaa ud fra en eneste Tanke. Al Mystik og romantisk
Spekulation rammes forsaavidt af denne Kritik. Dog
vil konsekvente Mystikere indrømme dens Resultat, idet
de hævder, at hvis Gud er det absolut Ene, kan der intet
siges om ham, jo, han er intet, d. v. s. intet af, hvad man
kan nævne eller tænke\footnote[1]{Smlgn. \textit{Religionsfilosofi.} § 21. Slutning.}.

Platon selv har banet Vejen for den rette Forstaaelse
af menneskelig Tænknings Virkemaade og Opgave, idet
% PRINTED AS IS: „«Staten«“, the opening guillemet turned, p. 32, seen at the image
han (i «Staten« og i »Faidros«, og senere i »Theaitetos«
og i »Sofisten«) lærer, at al Tænken er en Sammenfatten
gennem Vexelvirkning af Forbindelse og Sondring. Han er
derved Forgænger for den nyere Tids Synteseteori hos
Leibniz, Kant og deres Efterfølgere. ---

Den næste Betragtning (p. 142 B—155 E) gaar ud fra en
Forudsætning, som egentlig burde være afvist, naar man
stolede paa Resultatet af den første Betragtning. Thi som
vi saa, maatte Prædikatet »Væren« (οὐσία) afvises, hvis
der skulde gøres Alvor af »det Ene«. Men den nye Betragtning
gaar netop ud fra, at det Ene »er«. Den første
Betragtning gik ud fra at det, der er, er »Et« (εἰ ἕν ἐστιν);
den nye Betragtning gaar ud fra, at der kun kan tillægges
det Ene »Væren« (ἕν εἰ ἐστιν). Hist blev der gaaet ud fra
Begrebets Indhold; her gaas der ud fra, at dette Indhold
har Eksistens, --- at ἕν er et ὄν. Drøftelsen af denne Forudsætning
fører til modsatte Resultater af Drøftelsen af den
% --- p. 33 ---
\setcounter{footnote}{0}%
\opage{33}forrige. At Forudsætningen alligevel tages op og gennemgaas meget udførlig, er et Vidnesbyrd om det foreløbige og
prøvende i denne anden Del af Dialogen. Desuden er Drøftelsen
rigtignok ogsaa et Vidnesbyrd om, hvilket Sofisteri
Platon kunde komme ind paa, før hans Tanke havde vundet
fuld Klarhed.

Tankegangen er denne. Naar det Ene siges at være,
faar vi to Begreber, det Ene og dets Væren. Men disse
Begreber maa høre nøje sammen og udgøre en Helhed.
Enhver Helhed bestaar af Dele, og enhver af Delene maa
igen have Karakteren baade af Enhed og af Væren, og
disse to Prædikater maa igen danne en Helhed --- og saaledes
i det uendelige i Kraft af den engang indslaaede
Tankegang (κατὰ τὸν αὐτὸν λόγον ὡσαύτως). Men det er
ikke blot en Tohed, der saaledes gentager sig; det er en
Trehed, thi de to Prædikater og deres Forbindelse udgør
en Trehed, --- og ogsaa her kan vi fortsætte paa samme
Maade, saa at vi faar Brug for hele Talrækken. Det Ene
bliver da til en Mangfoldighed. Og som Helhed er det afsluttet,
har baade Begyndelse og Ende. Det er fremdeles
baade i Hvile og i Bevægelse, fremskridende gennem forskellige
Øjeblikke, paa hvert Punkt staaende mellem et
»nu« og et »derpaa«. --- Da der saaledes er Mulighed for en
hel Række hvert for sig paaviselige Prædikater, kan
Iagttagelse og Videnskab nu blive mulig, hvad de ikke
kunde efter den forrige Betragtning. ---

Selv Zeller og Jackson, der ellers véd at finde dyb
Filosofi i dette Afsnit af Dialogen, opgiver det dog her.
Der er ogsaa en aabenbar Sofisme i den Maade, hvorpaa
Begrebet Helhed her indføres. Det, at der skelnes mellem
to Begreber eller Prædikater (Enhed og Væren), viser paa
en Maade, at der er en Helhed, eftersom Prædikater, der
% --- p. 34 ---
\setcounter{footnote}{0}%
\opage{34}tillægges samme Emne, maa staa i en vis Sammenhæng
med hverandre. I Dialogens første Del har Platon jo ogsaa
hævdet, at der ingen Modsigelse er i, at samme Emne
hører ind under flere Ideer. Det er ganske vist en ubestemt
Helhed; det er den, der fremtræder i den tredie aristoteliske
Slutningsfigur (B er A og B er C), hvor Vanskeligheden
netop er den at finde, hvilken Sammenhæng man
har Ret til at antage mellem forskellige Prædikater til
samme Subjekt (mellem A og C). Det er i hvert Tilfælde
en rent logisk Helhed, en Tankesammenhæng, og det er
uberettiget at overføre paa den, hvad der gælder for Helheder
i Tid og Rum, Helheder, som er givne i virkelig
eller mulig Erfaring, --- reale Helheder. Og dog gaas der
uden videre ud fra, at den Helhed, Begreberne Enhed og
Væren udgør, maa underligge Rummets, Tidens og Tallets
Kategorier. Fra den historiske Parmenides’ Standpunkt, hvor
det »Ene« fra først af tænkes som en overalt ensartet
Kugle, er i hvert Tilfælde Rummets Kategori anvendelig;
men Anvendeligheder naas der ikke ved Slutning, men er
forudsat. For Platon er denne Tankegang helt fremmed,
afset fra den Begrebsforveksling, hvorved den begrundes. ---

Det er intet Tilfælde, at Tankegangen ud fra det »Ene«
som værende (eksisterende) faar en realistisk Karakter.
Naar der spørges, hvad vi mener med, at noget er, eksisterer,
viser allerede psykologisk Undersøgelse, at vi derved
støtter os til Iagttagelser, der godtgør, at dette noget staar
i fast og ubrydelig Sammenhæng med alt andet, der foreligger
for os. Tvivles der om noget af dette andet, rejses
blot det samme Spørgsmaal paany. At være (eksistere) er
at være Led i en lovmæssig Sammenhæng. Naar det altsaa
skal drøftes, hvad der ligger i, at det »Ene« er, maa man
nødvendigvis blive klar over, at man derved paa Forhaand
% --- p. 35 ---
\setcounter{footnote}{0}%
\opage{35}stiller det i Forhold til noget andet, og derved forstaas
forsaavidt, at Tankegangen glider over til Anvendelse af
Begreber som Rum, Tid og Tal, som netop er Former for
Sammenhæng. Man kommer ind i Relationernes Verden,
saasnart man spørger om Eksistens. Og selv om man holder
sig til det blotte Begrebsindhold, som i de rent formale
Videnskaber, undgaar man ikke Forholdet til noget
andet. Logisk og mathematisk er alt muligt, hvad der ikke
indeholder Selvmodsigelse; men selv dette formale Sandhedskriterium
forudsætter, at det »Ene« kan sættes i Forhold
til noget andet; ellers kan dets logiske Gyldighed
ikke begrundes.

Man kunde spørge, om det ikke er en for Platon fremmed
Betragtning, som her gøres gældende angaaende »Væren«
og Eksistens. Hertil maa svares, at Platon selv (i
tiende Bog af »Staten«) anvender det samme Kriterium
paa Væren, som her er gjort gældende. Det er i Anledning
af Forholdet mellem Kunst og Videnskab, at Platon kommer
ind paa dette Punkt. Kunsten holder sig til det enkelte
Billede, som det fremtræder, naar Genstanden ses fra
en enkelt, bestemt Side. Videnskaben naar til det sande
Begreb om Genstanden ved at støtte sig til Iagttagelser af
Genstanden fra saa mange forskellige Sider som muligt.
Et og samme Bord tager sig forskellig ud for os, alt efter
den Side, fra hvilken, og den Afstand, i hvilken vi ser den.
En Stok, der under sædvanlige Forhold ses som lige, vil,
naar den staar i Vand, ses som skæv. Under forskellige
Forhold kan en og samme Ting staa som blød og som
haard. Vil vi have Tingenes virkelige Beskaffenhed at vide,
maa vi tælle, maale og veje; vi kan da danne os et Begreb
om dem, ud fra hvilket vi under Hensyntagen til de
forskellige Omstændigheder kan forklare, hvorledes den
% --- p. 36 ---
\setcounter{footnote}{0}%
\opage{36}kommer til at tage sig forskellig ud. Væren tillægger vi det,
et saaledes dannet Begreb indeholder. (Staten X p. 598 A;
602C—603 A. Kfr. VII p. 528 AB). --- Vi vil siden faa Lejlighed
til at se, at Platon i sine senere Skrifter lægger stor
Vægt paa denne Betragtning, der klinger igen i Galilei’s
berømte Opfordring til at maale alt, hvad maaleligt er, og
gøre det maaleligt, som endnu ikke er det. Det er meget
mærkeligt, at Platon ikke har gjort Brug af den i »Parmenides«.
Den vilde have støttet den Kritik af Begrebet
det »Ene«, som han forud har givet, og vilde have ført til
det Resultat, der senere arbejder sig frem i »Filebos« og i
»Timaios«, nemlig, at det mystiske Begreb om det »Ene«
kun faar rationel Betydning ved at ses som Udtryk for
Emnernes formale og reale Sammenhæng. Spinoza’s »Substans«
er netop Parmenides’ »Ene« som Princip for Verdens
indre, lovmæssige Sammenhæng\footnote[1]{\textit{Spinoza’s Ethica} p. 9---11; 14---18. (Vidensk. Selsk. Skrifter 1918).}. ---

Dialogens sidste Stykke (p. 164 B—166 C) er Kontraprøve
paa den første Betragtning (der gik ud paa det Enes
Gyldighed). Det gaar ud fra, at det »Ene« ikke er, og
undersøger, hvad Følgen heraf bliver for det Mangfoldige
(τὰ πολλά). De mangfoldige Emner staar da blot i Forskelsforhold
til hverandre, og der vil aldeles ingen Sammenhæng
være mellem dem. Selvom de tilsyneladende danner
Mængder eller Masser (ὄγκοι), vil disse ved nærmere Betragtning
opløse sig i Kaos, i en Uendelighed af forskellige
Emner. Og ved ethvert tilsyneladende Hele vil man finde
stadig nye Dele indenfor det og stedse flere Yderled ved
dets Grænse. Tilsidst vil selve Mangfoldigheden forsvinde,
da dens Dele hver for sig skulde være Enere, men intet
Ene jo eksisterer! ---

% --- p. 37 ---
\setcounter{footnote}{0}%
\opage{37}Det kan her tydelig spores, at Meningen med »det Ene«
er, at det er Principet for lovmæssig Sammenhæng i Verden.
Falder dette Begreb fuldstændig bort, kommer vi til
en Grænse for vor Erkendelse, der er modsat den, som
den første Betragtning (137 A—142 A), der drøftede Indholdet
af »det Ene«, førte til. Ligesaa umuligt som det er at
fastholde Begrebet »det Ene« i dets fulde Konsekvens,
ligesaa umuligt er det at opgive det helt. Der er altsaa to
indbyrdes modsatte Forudsætninger, som hver for sig vilde
gøre Erkendelse umulig. Jeg har i min Fremstilling af Erkendelseslæren
kaldt den ene for »den absolute Identitetsrække«,
den anden for »den kaotiske Forskelsrække«\footnote[1]{\textit{Den menneskelige Tanke} p. 162---172. --- Se ogsaa \textit{Totalitet som Kategori} p. 8---11. (Vidensk. Selsk. Skr. 1917).}.
Mellem begge bevæger al menneskelig Erkendelse sig under
stadig Brydning af Identitet og Differens. ---

Dialogen ender pludselig efter den sidst omtalte Drøftelse,
uden noget Tilbageblik, uden afsluttende Bemærkninger,
uden nogensomhelst Antydning af den Interesse, de
forskellige drøftede Forudsætninger kunde have for det i
Dialogens første Del omhandlede Spørgsmaal, altsaa uden
at pege paa, hvad der kunde træde istedetfor den saa stærkt
kritiserede Idelære. Det taler stærkt for den Formodning,
at vi her kun har et løst Udkast. ---

Jeg har her kun omtalt anden Dels Begyndelse (p. 137 C
—155 E) og Slutning (p. 164 B—166 C). De Betragtninger,
som ligger imellem (p. 156 A—164 A), gaar jeg (afset fra et
enkelt Punkt) ikke ind paa. Jeg kan ikke finde noget
frugtbart ved dem, og Begrebsforvekslingerne florerer frodigere
der end i andre Dele af Dialogen. --- Det enkelte
Punkt, jeg vil dvæle ved, kræver et særligt Afsnit.
% --- p. 38 ---
\setcounter{footnote}{0}%
\opage{38}\chaphead{3. Begrebet det Pludselige (τὸ ἐξαίφνης) i Platons Filosofi.}

Undersøgelsen af »det Ene« som værende (142 B—155 E)
havde ført til en nærmere Udvikling af en Tanke, som
den platoniske Sokrates straks i Begyndelsen havde hævdet
(129 B), nemlig at flere Ideer meget vel paa én Gang
kunde gælde for samme Emne, en Erkendelse, som iøvrig
allerede var vunden i tidligere Dialoger (Fajdon 102 C;
103 B. --- Staten 479 A; 523 C). Og i et mærkeligt Stykke,
som afslutter den nævnte Undersøgelse, kommer Platon
(den platoniske Parmenides) nu ind paa Spørgsmaalet om
et successivt Forhold til de Ideer, der kan gælde for samme
Emne, --- om Overgangen fra Lighed til Ulighed, fra
Hvile til Bevægelse, fra Lidenhed til Storhed o. s. v. Og
her dukker Begrebet det Pludselige op, som frembyder en
særlig Interesse i Platon’s Filosofi.

a. Indførelsen af dette Begreb kommer selv meget
pludselig. Under Paavisningen af, at »det Ene« som
værende maa kunne have mange forskellige Prædikater,
fremhæves det særlig, at det maa være Tiden underlagt;
ellers kunde det ikke gaa over fra en Tilstand til en anden.
Men der lægges --- med en vis Tilslutning til Zenon
--- Vægt paa, at det i hvert Nu maa være, hvad det er, Nu’et
kan ikke forbigaas, uagtet der stadig maa skrides frem fra
Øjeblik til Øjeblik, saa at Nu’et straks maa slippes igen og
det næste Nu (τὸ ἔπειτα) gribes, (152 CDE). Her antydes
der endnu ikke noget alvorligt Problem. Men saa tages lidt
efter (155 E) Spørgsmaalet op igen, idet Modsætningen eller
Modsigelsen mellem de forskellige Prædikater (eller de Tilstande,
disse angiver) skarpt fremhæves. Der maa ske en
Forandring (μεταβολή), en Overgang fra en Tilstand til en
helt anden Tilstand. Den ene Tilstand maa slippes, den
% --- p. 39 ---
\setcounter{footnote}{0}%
\opage{39}anden overtages. Men naar sker denne Overgang? Saalænge
den tidligere Tilstand (Lighed, Hvile, Lidenhed o. s. v.)
er der, sker den ikke, og naar den nye Tilstand (Forskel,
Bevægelse, Storhed o. s. v.) er der, er Overgangen allerede
sket. Den sker pludselig. Men dette Pludselige (τὸ ἐξαίφνης)
er et Paradox (ἄτοπον), thi det betyder en Overgang mellem
to Tilstande, og da der ikke kan gives nogen Tid, i
hvilken et Emne hverken er i Hvile (o. s. v.) eller i Bevægelse
(o. s. v.), saa er det Pludselige ikke i nogen Tid.
Til dette Pludselige gaar den gamle Tilstand over, og af
dette dukker den nye Tilstand op.

Den Løsning, som her antydes, bestaar i at anerkende
Overgangens, og dermed det Pludseliges Realitet, men at
sætte den udenfor Tiden. Hvad der er udenfor Tiden og
dog har Gyldighed, det maa ifølge Platon’s Filosofi høre
til Ideverdenen, og Platon maatte konsekvent have antaget
en Pludselighedens Ide. Det vilde han sikkert endnu nødigere
have gjort end at antage en Ide for Smudset. Thi denne
nye Ide vilde bringe Uro ind i Ideernes rolige Verden. I
dobbelt Forstand bliver det Pludselige da noget, der ingensteds
har hjemme (et ἄτοπον); hverken Ideernes eller Emnernes
Verden kan huse det. Men ved nærmere Eftersyn
berettiger Platon’s Tænkning ham ikke til at sætte det
Pludselige udenfor Tiden. Det udfylder, skønt det ikke
skal være Tilstand, dog et Led i Tilstandenes Række, er
altsaa et Emne (et Fænomen, noget, der foreligger). Begrundelsen
af dets Tidløshed beroede jo kun paa, at Overgangen
hverken skulde være Hvile (o. s. v.) eller Bevægelse
(o. s. v.). Men Fejlen hos den platoniske Parmenides kunde
jo netop ligge i, at der forudsættes et kontradiktorisk Forhold
mellem disse Begreber. Naar samme Emne godt kan
falde ind under forskellige Ideer, er et Baade—Og ikke
% --- p. 40 ---
\setcounter{footnote}{0}%
\opage{40}paa Forhaand udelukket; Lighed og Forskel, Hvile og Bevægelse
o. s. v. behøver ikke at være indbyrdes modsigende
Begreber. Det er ikke nødvendigt, at en Ide forsvinder,
for at en anden kan komme til at gælde. Overgangen kan
altsaa ske fra en Tilstand af Lighed med underordnede
Forskelligheder til en Tilstand af Forskellighed med underordnet
Lighed, --- fra en urolig Hvile til en ensformig Bevægelse
o. s. v. I nyere Tid gik endnu Descartes ud fra et
kontradiktorisk Indhold mellem Hvile og Bevægelse, men
Leibniz viste Muligheden af en Overgang fra Hvile til Bevægelse,
naar man opfattede Hvilen som en Ligevægtstilstand,
i hvilken forskellige Bevægelsestendenser holdt hinanden
Stangen. Men ligesom Spørgsmaalet »Væren« (Eksistens)
ikke kunde besvares ved at holde sig til et enkelt Væsen
for sig alene, men kun ved at se det i Forhold til andre
Væsener, saaledes kan Overgangen fra en Tilstand til en
anden ikke forstaas ud fra den enkelte Tilstand som isoleret;
kun under visse bestemte Betingelser, navnlig under
udefra kommende udløsende Paavirkningers Indflydelse sker
Overgangen. Det er Isolationen, det gælder at ophæve,
baade indenfor Ideernes og indenfor Emnernes Verden. ---

Hos to moderne Tænkere har Platon’s Lære om det
Pludselige vakt Interesse.

Kierkegaard, der i en Note til »Begrebet Angst« (Kap. 3)
drøfter denne Lære, oversætter »det Pludselige« ved »Øjeblikket«
og siger, at det for Platon er »det Ikkeværende
under Tidens Bestemmelse«. Dette vilde Platon ganske sikkert
ikke indrømme. Han kender ingen tom Tid. Han
paastaar jo udtrykkelig, at det Pludselige ikke er i nogen
Tid. Men Kierkegaard har forsaavidt Ret, som Platon ikke
konsekvent kan sætte det Pludselige udenfor Tiden, da det,
hvortil en Tilstand gaar over, og hvorfra en anden Til%
% --- p. 41 ---
\setcounter{footnote}{0}%
\opage{41}stand begynder, selv maa være Led i Tidsrækken. Naar
Kierkegaard uden videre oversætter (τὸ ἐξαίφνης) ved »Øjeblikket«,
har han derved allerede korrigeret Platon; thi
ethvert Øjeblik er et Led i Tidsrækken. Desuden havde
Kierkegaard Brug for en positiv Betydning og en bestemt
Tidsbeliggenhed for det »Pludselige« i sin Lære om »Springet«.
Springet, den absolute Frihedsakt skal være en positiv
Handling i et bestemt Øjeblik og skal kunne iagttages.
Men det skal ikke som andre Tilstande eller Øjeblikke
være lovbestemt; det er ikke visse bestemte Udløsningsbetingelser
underlagt. »Det Pludselige«, siger Kierkegaard
(Begrebet Angst, p. 137), »kender ingen Lov«. Som jeg har
paavist i min Bog om Kierkegaard som Filosof (anden
Udgave p. 78—82) henlægger Kierkegaard alligevel »Springet«
mellem to Øjeblikke; i det første er det endnu ikke
sket, i det andet er det allerede sket. Herved nærmer han
sig igen til Platon’s egen Opfattelse af det Pludselige. Men
saa kan det Pludselige ikke kaldes »Øjeblikket«, naar man
ved et Øjeblik forstaar et selvstændigt Led i Tidsrækken.
De to Øjeblikke, mellem hvilke Springet ifølge Kierkegaard
finder Sted, har intet Øjeblik imellem sig. Mærkværdigt
nok mener Kierkegaard, at Springet alligevel skal kunne
iagttages.

Natorp har set, at ifølge Platon er det Pludselige ikke
i Tiden, og han finder heri en Støtte for sin Fortolkning
af Platon’s Idelære i Analogi med den moderne Erkendelsesteori
(særlig Marburgerskolens). Han finder nemlig deri
den Opfattelse udtrykt, at ikke Sansningen, der er bunden
til Tiden som diskontinuerlig Succession, men kun Tanken
kan hævde Kontinuitet, begrunde Overgangen fra et Prædikat
til et andet. Medens Kierkegaard fortolker det Pludselige
som Springet, sætter Natorp uden videre »Overgangens
% --- p. 42 ---
\setcounter{footnote}{0}%
\opage{42}Kontinuitet« som ensbetydende med »Overgang«; at Overgangen
sker »im Auf-einmal« (Parm. 156 D) skal betyde,
at »Overgangens Kontinuitet har sin Rod i en Tænken«\footnote[1]{Afsnittet »Platon« i Aster’s »Grosse Denker« I, p. 127.}.
Men Platon betragter tydelig nok den pludselige Overgang
som et Faktum, det gælder at forstaa, og hans Begrundelse
af, at det ikke kan være i nogen Tid, er rent negativ, uden
nogen Henvisning til Tankens, Ideernes Verden. Det vilde
heller ikke nytte at henvise Overgangen til den rene Logik,
dels fordi den der bliver til kontinuerlig Overgang, medens
Gaaden netop er det diskontinuerlige i Overgangen fra Ide
til Ide, dels fordi der saa kommer et nyt Problem frem,
nemlig hvorledes Forholdet er mellem den logiske Kontinuitet
(der egentlig er Identitet) og den reale (temporale)
Diskontinuitet. Men dette er nu rigtignok netop et Spørgsmaal,
hvis Betydning Marburgerskolen overser\footnote[2]{Smlg. \textit{Totalitet som Kategori} p. 46---48. (Vid. Selsk. Skr. 1917).}. Denne
Skole staar her overfor samme Vanskelighed som Parmenides
og Platon, en absolut Modsætning mellem Kontinuitet
og Diskontinuitet, og den undervurderer Diskontinuitetens
stadige Betydning i vor Erkendelse som ansporende
og problemstillende Faktor.

Diskontinuitetens Kendsgerning, den Pludselighed, hvormed
en Overgang kan ske, hænger nøje sammen med, at
ethvert Emne staar i Forhold til andre Emner. Hvis vi
kun raadede over en eneste Oplevelse, vilde Begrebet
Diskontinuitet ikke være dannet. Et enkelt Punkt for sig
alene kan ikke kaldes diskontinuerligt. Og Pludselighed kan
ikke udsiges om nogen Oplevelse, der ikke staar i kvalitativ
Modsætning til andre Oplevelser. Nærmere Undersøgelse
viser, at vor Erkendelse udvikler sig under stadig fornyet
Brydning mellem Kontinuitet og Diskontinuitet. Denne
% --- p. 43 ---
\setcounter{footnote}{0}%
\opage{43}Overbevisning nærmede Platon sig, da han satte en Verden
af forskellige Ideer istedetfor Parmenides’ »Ene«. Dermed
var ikke alle Vanskeligheder overvundne, men der var
tilvejebragt bedre Betingelser for at kæmpe med dem. Og
trods alle Uklarheder i Dialogen »Parmenides« er Platon
dog her bleven opmærksom paa, at det ikke er nok at
skelne mellem Ideer og Emner, men at Problemet Enhed
—Mangfoldighed ogsaa stiller sig indenfor Idelæren. Det
Spørgsmaal, hvorledes et Emne fra at høre ind under en
Ide kan komme til at høre under en anden Ide, fører Platon
til paa sin Vis at tage Zenon’s Problem op. At han
anerkender det som et virkeligt Problem, medens Parmenides
og Zenon betragter det som Frugt af en Illusion,
kommer af, at Ideverdenen for ham er en Mangfoldighed.
Hver enkelt Ide er for ham en Realitet; derfor maatte
Overgangen fra en Ide til en anden ogsaa være en Realitet.
Zenon lagde kun i polemisk Hensigt Forudsætningen
om en real Mangfoldighed til Grund. Platon er nu, som
vi har set, enig med Zenon deri, at ud fra en absolut
Mangfoldighed (en kaotisk Forskelsrække) er ingen Erkendelse
mulig (Parm. p. 164 B—166 E). Men ligefra sin Filosoferens
Begyndelse havde Platon, i Tilslutning til Herakleitos,
Tanken om, at det i Erfaring givne, Emnerne, var
stedt i stadig Vorden og Forandring. Da der nu til hver
Gruppe af de i Forandring stedte Emner svarer en Ide,
maa Problemet om Emners Overgang fra en Ides Omraade
til en anden Ides Omraade nødvendigvis komme til at
fremstille sig med en ganske anden Styrke, end naar man
blot helt hypotetisk taler om Bevægelse og Forandring. Og
her springer da det Pludseliges Problem frem med Kraft. ---

b. Begrebet det Pludselige er Platon her stødt paa ved
Modsætningen mellem de skarpt sondrede Ideer og Emner%
% --- p. 44 ---
\setcounter{footnote}{0}%
\opage{44}nes kontinuerlige Strøm. Men ogsaa paa et helt andet Sted
i Platon’s Filosofi er dette Begreb dukket frem. Medens
Parmenides og Zenon ikke var Psykologer, er psykologisk
Interesse et meget fremtrædende Træk hos Platon. Det var
ikke blot Dialektik og Metafysik, der førte ham ud over
det sokratiske Standpunkt; det var nok saa meget hans
dybe psykologiske Forstaaelse. Hans Lære om Eros (i
Symposion og Faidros) er med Rette bleven kaldt den
antike Tænknings vigtigste Bidrag til Psykologien. Medens
Sokrates en Gang for alle forudsatte, at der hos Mennesker
maatte være baade Trang og Evne til at danne Tanker,
der kunde fastholdes under alle Omskiftelser, og kun bestræbte
sig for at faa dem til at lade Trangen til Klarhed
og Sandhed faa Indflydelse paa deres Liv, søger Platon at
vise, hvorledes denne Trang, den intellektuelle Higen, og
den tilsvarende Evne kan udvikle sig Skridt for Skridt.
Eros, der for ham i videste Betydning betegner hele Driftslivet,
optræder først som Trang til sanselig Tilfredsstillelse,
men kan dog pege ud over denne. Og selve den Trang,
der rører sig efter sanselig Tilfredsstillelse, kan vise sig at
tjene Maal, der ligger ud over de enkelte Individer hvert
for sig. Gennem Forplantningstrangen virker det enkelte
Individ uden at vide det for Slægtens Bestaaen, og i Forældreinstinktet
viser det sig ganske tydelig, at her rører
sig noget, som fører ud over individuel Selvopholdelse. Den
Evne og Trang til at skabe og sikre Liv, der ytrer sig i
disse Drifter, vækkes af skønne, tiltalende Genstande, der
kommer til at staa som Formaal fra Individernes Trang.
Men fra legemlig Skønhed kan der gaas over til aandelig
Skønhed som udløsende Kraft, idet der er et Ligheds- eller
Slægtskabsforhold (ἀδελφόν, συγγενές, Symp. 210 B C) mellem
alt, hvad der er skønt. Og i Kraft af samme Forhold
% --- p. 45 ---
\setcounter{footnote}{0}%
\opage{45}kan der gaas over fra praktisk Stræben til intellektuel
Stræben, Stræben efter Sandheden for dens egen Skyld.
Gennem en Række af Trin kan Individet saaledes hæve
sig til Erkendelse eller rettere Skuen af de højeste Ideer. ---

Den Tankegang, som ligger til Grund ved denne platoniske
Opfattelse, har sin Betydning den Dag idag. Naar
for Platon Elskovstrang, Ærestrang, digterisk Trang og
Erkendelsestrang kan danne Led i en Udviklingsrække, saa
er det, fordi de kan danne en Række af, hvad man kan
kalde psykologiske Ækvivalenter. Det er en Række af
Motivforskydninger, han beskriver i sin Lære om den store
Eros. Der er Idealitet i enhver Drift, forsaavidt den peger
ud over den givne Tilstand, ud over hvad der hidtil er
opnaaet, --- stiller en Fordring. Og heri ser Platon Muligheden
af en Metamorfose, idet en Formaalstanke formedelst
Lighed kan afløse en anden. Eros kan blive den store
Læremester, idet Driften, naar den engang er vakt, kan
faa nye, maaske højere Formaal\footnote[1]{Smlgn. min Afhandling om Symposion i tredie Bind af Mindre Arbejder.}. Hvis dette er rigtigt,
skjuler der sig en dyb psykologisk Indsigt i Platon’s Lære
om Eros, og den er andet og mere end »ein geistreiches
Spiel«, som Wilamowitz-Moellendorff nylig har kaldt den\footnote[2]{\textit{Platon,} II, p. 172---174.}.

Men hvad man saa end mener herom, er det aabenbart,
at Platon paa Driftslivets Omraade antager Muligheden
af en psykologisk Udviklingsproces, indenfor hvilken
der i Kraft af bestemte Love sker Overgang mellem højst
forskellige Stadier. Livets Øjeblikke staar her ikke, som
de efter Zenon’s Forudsætninger maatte, med absolut Selvstændighed
overfor hverandre, men danner (eller kan dog
danne) en Række, i hvilken de senere Led stedse forud%
% --- p. 46 ---
\setcounter{footnote}{0}%
\opage{46}sætter de foregaaende, saa at Leddene ikke kan byttes om;
Successionsordenen er bestemt; --- Addendernes Orden er
ikke ligegyldig. Paa ethvert Punkt er Individet i Hvile,
forsaavidt dets Trang, saaledes som den har udviklet sig
gennem de forudgaaende Stadier, finder Fyldestgørelse;
men det er tillige i Bevægelse, forsaavidt det opnaaede
Trin (ἐπαναβαθμός) danner Forudsætning for eller indeholder
Tilskyndelse til at gaa videre. Kun naar der skrides
frem i en bestemt Rækkefølge (ἐφεξῆς καὶ ὀρθῶς), naas
det højeste Maal, et Menneske kan naa. Kun den naar
dette Maal, der ved at gaa gennem Stadiernes Række er
bleven styrket og er vokset (ῥωσθεὶς καὶ αὖξηθείς) (Symp.
210 D E). Men saa kan tilsidst den højeste Skønhed pludselig
(ἐξαίφνης) aabenbare sig for ham (210 E).

Medens det Pludselige optraadte pludselig og overraskende
fra Eleatismens og Idelærens Standpunkt, hvor man gik fra
Væren til Væren, fra Ide til Ide, kommer det Pludselige
paa en Maade ikke pludselig i »Symposion’s« Forskydningslære,
der indskærper den bestemte Rækkefølge af de
forudgaaende Stadier, som afsluttes med den højeste Skuen.
Der gaar en psykologisk motiveret Forberedelse forud for
det sidste Led i Rækken, som dog ogsaa hører med til
Rækken. Det er en Art Kontrastvirkning, der gør sig gældende,
og en saadan har det mærkelige, at den baade adskiller,
nemlig ved Kvalitetsmodsætningen mellem Leddene,
og forener, nemlig ved at det følgende Leds Kvalitet er
bestemt netop ved Modsætningen til det foregaaende. Overgangen
fra næstsidste til sidste Led er bestemt ved den Lov,
som gælder for Overgangen mellem alle Led i Rækken.
Selvom det sidste Led er nok saa værdifuldt i Forhold til
de foregaaende, --- selv om det betegner en Tilstand, hvor
Livet først ret faar Værdi for Mennesket, --- saa betegner
% --- p. 47 ---
\setcounter{footnote}{0}%
\opage{47}det dog et bestemt Tidspunkt, der har sin bestemte Plads
indenfor Livets Gang. (Symp. p. 211 D: ἐνταῦθα τοῦ βίου,
ἐίπερ που ἄλλοθι, βιωτὸν ἀνθρώπῳ). Det pludselig Optrædende
kommer som en moden Frugt, betinget ved Fortidens
Vækst og Blomstring.

I sin interessante Studie »La théorie platonicienne de
l’amour« (Paris 1908) fremhæver Léon Robin stærkt den
komplexe, syntetiske Karakter af den platoniske Eros. Der
gør sig i denne Følelse eller Drift kontrære Modsætninger
gældende, og Robin mener med Rette, at der paa dette
Punkt i Platon’s Tænkning viser sig en Indflydelse af
Herakleitos; denne Filosof lagde jo saa stor Vægt paa de
Kontraster, Tilværelsen frembyder. Og deri, at der gennem
den paa Modsætninger rige Tilstand, Eros betegner, kan
naas et inderligt Forhold til Ideerne, finder Robin en Tilnærmelse
til at mildne den skarpe Forskel mellem Ideer
og Emner.

Da jeg i denne Afhandling er kommen ind paa Læren
om Eros i Anledning af Begrebet det Pludselige, er det
ganske naturlig de successive Kontraster i den af Eros udspringende
psykologiske Proces, der interesserer mig. Det
siges jo udtrykkelig, at Eros som Søn af Poros og Penia
snart er rig, snart fattig, snart livskraftig, snart døende, og
det er gennem saadanne modsatte Faser, at Mennesket
vokser til at opnaa det store Øjeblik, da det Højeste pludselig
aabenbares. Det er saa at sige gennem en Kamp for
Tilværelsen, at der naas til dette Højdepunkt.

Deri, at der gennem psykologisk Udvikling kan skrides
frem fra Delagtighed i en Ide til Delagtighed i en anden
Ide, ligger allerede en Antydning af, at Ideerne indbyrdes
maa staa i et indre Forhold til hinanden. At de samtidig
kan være forenede i samme Emne, pegede allerede i denne
% --- p. 48 ---
\setcounter{footnote}{0}%
\opage{48}Retning. Men det fremtræder endnu tydeligere ved, at de
(eller Forholdene til dem) kan betegne forskellige Stadier i
en kontinuerlig Udvikling. Selv om de her kun mødes for
straks at skilles, saa hævder jo Platon selv i sin Dialektik,
at det baade er Tankernes Forskellighed og deres Lighed,
der bringer dem i Forhold til hverandre. Platon’s Idelære
og hans Lære om Eros peger i hver sin Retning, idet Eros
er uforstaaelig ud fra Idelæren, der ikke indeholder nogen
Mulighed for at tillægge Stræben og Udvikling nogen
virkelig Betydning. Men som jeg har søgt at vise, indeholder
Læren om Eros alligevel Synspunkter, der kan føre
til at raade Bod paa Idelærens Vanskeligheder. Det er Erfaringserkendelsens
Modsætning til Spekulationen, der her
gør sig gældende. Og efterat Platon i »Parmenides« saa
skarpt havde fremhævet disse Vanskeligheder, var det at
vente, at han ogsaa ad anden Vej end den psykologiske
Vej, han digterisk havde fulgt i »Symposion« og »Faidros«,
vilde tage de Problemer op, som her maatte stille sig. ---
Inden jeg gaar over til at undersøge de nye Tankeforsøg,
Platon saaledes kommer ind paa, vil jeg først fremdrage
nogle Antydninger, der fører i lignende Retning som Læren
om Eros, idet de skildrer en aandelig Udvikling henimod
Ideopfattelsen som Maal.

Ikke blot i »Symposion«, ogsaa i »Faidon« og i »Faidros«
omtales det Pludselige i den højeste aandelige Tilstands
Optræden. Fra den lavere Verden, i hvilken vi
Mennesker lever, er det dog, ifølge »Faidon« (p. 109 E),
engang imellem muligt at dukke op og faa et Glimt af
Ideverdenen, ligesom Fiskene engang imellem kan dukke
op af Havet og ser noget af vor Verden. I »Faidros«
(p. 248 A) faar Vognstyreren paa den menneskelige Sjælevogn
under Farten i Præexistensen lige netop Hovedet op
% --- p. 49 ---
\setcounter{footnote}{0}%
\opage{49}over sin vante Horizont, saa han kan faa et Glimt af den
sande Virkelighed. Paa disse Steder udtrykkes i billedlig
Form Modsætningen mellem det højeste Trin af menneskelig
Eksistens og de forudgaaende Trin, som dog er de nødvendige
Forudsætninger. Det er, som Læren om Eros viser,
ikke en rent teoretisk Udvikling, der naar sit Maal paa
hint højeste Trin. Der forudsættes Udvikling af hele Personligheden,
en Tilegnelse ad Livets Vej. Allerede i »Symposion«
udtales dette, naar det siges, at der ikke blot finder en
Skuen Sted, men ogsaa Samliv med det Højeste (θεᾶσθαι
καὶ συνεῖναι p. 212 A). Og endnu bestemtere udtales det i
»Statens« syvende Bog og i det syvende Brev. Hist læres,
at det Højeste ikke er Erkendelse, men Forhold til noget,
der er hævet over al Erkendelse, og som vi derfor heller
ikke kan tillægge Eksistens i Ordets sædvanlige Betydning.
Skønt det ifølge Platon kun er muligt for den, der har
gennemgaaet en streng videnskabelig Skole, at naa dette
Punkt, saa er det dog en personlig Oplevelse, der tilsidst
er det afgørende, saa at kun indirekte og billedlig Meddelelse
til andre er mulig, og da først, efter at man ikke
blot med Tanken, men med hele Sjælen er bleven vendt i
Retning af det Højeste (Staten VII p. 518 C D). I det syvende
Brev (p. 341 C; 344 B) fremhæves skarpt baade det
Pludselige og den personlige Oplevelse. Som en Protest
imod, at den unge Dionysios gjorde sig til af at være indviet
i Platons dybeste og inderligste Tanker, udtaler denne
her: »Dersom nogen nu eller senere paastaar at besidde en
Viden om det, som min Stræben egentlig gælder, hvad
enten de nu skulde have hørt det af mig eller af andre,
eller selv have udfundet det, saa erklærer jeg, at det er
umuligt, at de kan vide noget derom. Jeg har intet skrevet
derom og vil aldrig komme til at skrive noget derom, thi
% --- p. 50 ---
\setcounter{footnote}{0}%
\opage{50}det kan ikke meddeles som andre Erkendelser. Men efter
Arbejde dermed og efter at man har taget det op i sit Liv,
opblusser der pludselig (ἐξαίφνης) en Ild, som om en
Gnist var falden ind i Sjælen. Naar denne Ild først er tændt,
nærer den sig selv«\footnote[1]{Hvad der i denne Udtalelse fremhæves, det Pludselige og det Uudsigelige eller kun indirekte Meddelelige, det indskærper \textsc{Plotinos} ogsaa. Se \textit{Ennead.} V, 3, 17 (τότε δὲ χρὴ ἑωρακέναι πιστεύειν, ὅταν ἥ ψυχὴ ἐξαίφνης φῶς λάβῃ) og VI, 9, 4 (--- κατὰ παρουσίαν ἐπιστήμης) κρείττονα . . . οὖδὲ ῥητὸν οὖδὲ γραπτόν).}.

Ogsaa her er den pludselige Optræden af den nye Tilstand
et væsentligt Træk, samtidig med, at aandelig Vækst og personlig
Indleven er Betingelser for den. --- Psykologisk set udelukker
pludselig Opleven og lang Forberedelse ikke hinanden.

I Læren om Eros lægges der Vægt paa, at alle Sjælekræfter
virker sammen paa koncentreret Maade. Men selv
hvor Platon overvejende beskæftiger sig med den rent intellektuelle
Udvikling, er der for ham ingen Modsigelse i, at det
højeste Stadium, trods sin Modsætning til de foregaaende
Stadier, dog er forberedt ved den.

I »Faidon« (p. 74—79) skildres den strengt videnskabelige
Erkendelses Udvikling fra den Sammenhæng mellem
Forestillinger, der tilvejebringes ved Forestillingsassociation
formedelst Lighed og ved uvilkaarlig Sammenligning. Sanselige
Ting, som Træstykker og Stene, ligner hverandre i
visse Ting, ikke i andre, og fuldkommen Lighed (Identitet)
forekommer ikke i sanselig Erfaring. Og dog er det ved at
sysle med hine ufuldkomne Ligheder, at vi arbejder os
frem til Lighedens Ide! Og som det gaar med Lighed, gaar
det med andre Tanker, med Tankerne om Retfærdighed,
Skønhed og Godhed. I det, Erfaringen viser os, er der i
det højeste en Tendens, en Stræben efter, hvad Ideen indeholder,
men alligevel kan Erfaringen tjene som Forberedelse
% --- p. 51 ---
\setcounter{footnote}{0}%
\opage{51}til den fuldendte Erhvervelse af Ideindholdet, som kun den
højeste Videnskab skænker. At dette læres i »Faidon«, er saameget
interessantere, som denne Dialog ellers træder i en vis
Modsætning til »Symposion« ved at betone, hvor store
Hindringer for sand Erkendelse Sansning og sanselig
Drift kan være. Erfaring indeholder for Platon aldrig den
fulde Forklaring af den sande Videnskab. Han kender ikke
Induktion i moderne Betydning af Ordet, og bruger heller
ikke det tilsvarende græske Ord (ἐπαγωγή), der først forekommer
hos Aristoteles. Hvad der foregaar under den intellektuelle
Udvikling, er for Platon ikke en Indførelse i
Emnernes Sammenhæng, men »en Opløftelse (ἐπαναγωγή)
af det Bedste i vor Sjæl til Skuen af det Bedste i det sande
Værende« (Staten VII p. 532 C). Det er, som man træffende
har udtrykt det, en Art Himmelfart, Platon her beskriver\footnote[1]{\textsc{Teichmüller}: \textit{Studien zur Geschichte der Begriffe,} p. 407; 425. Teichmüller formoder, at der i den platoniske Epanagoge (Opløftelse) ligger Spiren til det aristoteliske Epagoge (Induktion), der er en Indførelse i det Almenes Verden. Himmelfarten bliver hos Aristoteles til en Indførelse i Emnernes almindelige Love.}.

I mere metodisk Form har Platon i Statens sjette Bog
skildret Erkendelsens Udvikling gennem en Række af Trin.
Fra Sansningen skrides der op til Mening og Formodning,
derfra til den strenge matematiske Erkendelse og tilsidst
til den Prøvelse af denne Erkendelses Forudsætninger, som
ifølge Platon fører ud over alt Hypotetisk til Erkendelse af
det i sig Værende. Der foregaar her en Fremskriden fra
Trin til Trin, hvor intet Trin kan undværes, hvis man vil
naa det højeste. ---

Platon har altsaa under forskellige Former fastholdt, at
saa højt hævet over alle Sansninger, Drifter, Hypoteser og
Slutninger den højeste Erkendelse ogsaa ligger, baade hvad
Klarhed og hvad Værdi angaar, saa naas den dog kun
% --- p. 52 ---
\setcounter{footnote}{0}%
\opage{52}gennem trofast Stræben og Arbejde paa de lavere Trin.
Men naar Spørgsmaalet opkastes om, hvorledes en saadan
Udviklingsgang, en saadan Opstigen overhovedet er mulig,
da staar den ganske vist for Platon som et Under, der
kun kan forklares ved, at vore Sjæle allerede i Præeksistensen
har (selv om det kun var glimtvis) staaet Ansigt
til Ansigt med det, der nu maa erhverves under Savn og
Længsel og gennem alvorligt Arbejde. Denne mytologiske
Forklaring er udsprungen af Forundringen over, at det
Resultat, der naas ad de lange Veje, synes langt at overgaa
hvad der forberedtes undervejs. Platon’s Udvej her kan
minde om den moderne Udviklingsfilosofis Løsning af Erkendelsesproblemet,
kun at man her peger tilbage fra de
enkelte Individers Erfaringer til Slægtens Erfaringer, medens
Platon peger tilbage til selve de enkelte Individers Præeksistens\footnote[1]{Smlgn. min Afhandling: \textit{Nogle Bemærkninger om Platon’s Psykologi.} (Tidsskrift for Filologi og Pædagogik, 1879), p. 209---213.}.
---

Inden jeg slutter dette Afsnit, hvor Læren om Eros har
spillet en væsentlig Rolle, maa det bemærkes, at denne
Lære ikke udtømmer hele Platon’s Psykologi. Der maa,
som jeg allerede udviklede i min Doktordisputats »Den antike
Opfattelse af Menneskets Villie« (1870) skelnes tre
psykologiske Typer hos Platon. Den ene er den, jeg i dette
Afsnit særlig har fremdraget, og som raader i »Symposion«
og i »Faidros«. Den er mærkelig ved at tage Udgangspunktet
i elementær Naturtrang og dog føre op imod et
absolut Højdepunkt. Bevægelsen foregaar ud ad én Linie.
Den anden psykologiske Type, Platon skildrer, frembyder
en Udvikling fra forskellige Udgangspunkter, idet der antages
en Forskel, henholdsvis en Modstrid mellem tre forskellige
Tendenser i Sjælen, --- til Nydelse, til Selvhævdelse
% --- p. 53 ---
\setcounter{footnote}{0}%
\opage{53}i Mod og Harme og til Erkendelse. Hvad der i Eros vældede
frem med samlet Kraft, virker her sporadisk, og Opgaven
bliver at naa en Harmoni mellem de stridende Tendenser.
I Statens fjerde Bog lægges denne Opfattelse til
Grund; den spores dog ogsaa i »Faidros«. I den tredie
Type forudsættes al Energi koncentreret i Tankelivet, som
raader over en Magt, der holder Nydelses- og Selvhævdelsestrang
nede. I »Faidon« raader denne Type; men den
optræder ogsaa i »Statens« niende Bog.

Forskellen mellem de tre Typer kan, som det især ses
af Dr. Hans Ræders Undersøgelser, spille en Rolle ved
Drøftelse af Dialogernes Kronologi. Men de gaar gennem
forskellige Metamorfoser over i hverandre. Det er et Vidnesbyrd
paa Rigdommen i Platon’s psykologiske Opfattelse,
at han har fremdraget dem alle; men deres Modsætning
stiller store Problemer. Og hvad der i denne Sammenhæng
især interesserer os, er, hvad jeg allerede i det nævnte Skrift
lagde Vægt paa\footnote[1]{\textit{Den antike Opfattelse af Menneskets Villie,} p. 64; 74.}, at hverken den Higen, som kaldes
Eros, eller den Udvikling til Harmoni, som den anden Type
betinger, kan udledes fra selve Idelæren. Idelæren forudsætter
et Sjæleliv, der en Gang for alle er paa Højden.
Hvorledes det er kommet derop, kan den ikke forklare, ---
og, som vi har set, og som Parmenides’ tredie Indvending
slaar fast, heller ikke, hvorledes den Erkendelse, der kan
vindes deroppe, kan faa virkelig Anvendelse paa Emnernes
Verden.

\chaphead{4. Senere Dialoger i Forhold til Problemerne i\\ „Parmenides“.}


Den Selvkritik, som Platon, ifølge den i det foregaaende
hævdede Opfattelse, har øvet i »Parmenides«, har afsat
% --- p. 54 ---
\setcounter{footnote}{0}%
\opage{54}Virkninger i hans senere Tænkning, skønt han ikke har
gendrevet de i »Parmenides« fremsatte Indvendinger mod
Idelæren. Hans senere Lære kan ikke betragtes som en
ligefrem Fortsættelse af Læren i »Symposion«, »Faidon«
og »Staten«, og dog har han ingenlunde opgivet Idelæren,
skønt den træder noget i Baggrunden. Han maatte, for at
kunne »gøre noget ved Filosofien« (Parm. 135 C), altsaa for
at hans Tænkning ikke skulde gaa i staa, anlægge en anden
Problemstilling end den, han hidtil havde benyttet. En
saadan ny Problemstilling spores i den hele Gruppe af
Dialoger, der ligger mellem »Parmenides« og »Lovene«,
hvilken sidste Dialog ikke behandler Spørgsmaal, der angaar
Idelæren. Denne Gruppe bestaar af »Theaitetos«, »Sofisten«,
»Statsmanden«, »Filebos« og »Timaios«. Nyere Forskere
(Jackson, Brochard, Ræder, Robin) er enige i, at
der er en indre Sammenhæng mellem disse Dialoger
indbyrdes, hvad Problemer og Tankegang angaar, selv om
det kronologiske Forhold mellem dem opfattes forskellig.

Jeg kan ikke følge Dr. Ræder i at sætte »Theaitetos«
før »Parmenides«, da den Maade, paa hvilken den tager
Erkendelsesproblemet op, kun er forstaaelig paa Baggrund
% defect: „indgaaende“, its second e broken (ę) in the print, p. 54; the layer read ç
af »Parmenides« indgaaende Kritik af Idelæren, og da Problembehandlingen
i »Theaitetos«, som jeg vil forsøge at
vise, er et Skridt i den Retning, der føres videre i »Sofisten«
og »Statsmanden« og afsluttes i »Filebos« og »Timaios«.
Derfor kan der godt, hvad sproglige Grunde menes
at tale for, være nogen Tidsafstand mellem »Theaitetos«
og »Sofisten«.

Allerede i »Staten« og i »Faidros« var Erkendelse beskreven
som Vekselvirkning af Forbindelse og Sondring,
Samskuen (Synopsis) og Analyse, og i denne Vekselvirkning
lagde sig for Dagen, hvad Platon kaldte Dialektik.
% --- p. 55 ---
\setcounter{footnote}{0}%
\opage{55}Heraf ses, at de abstrakte Udviklinger om Enhed og Mangfoldighed
i andet Afsnit af »Parmenides« aldeles ikke havde
været nødvendige for, at Platon havde kunnet fortsætte sin
Filosoferen, efterat han havde anerkendt Indvendingerne
mod Idelæren i dens oprindelige Form. Selve det Begreb
om Tænkningen som stadig Vekselvirkning af Forbindelse
og Sondring, som han i de to nævnte Dialoger var naaet
til, vilde have været et tilstrækkeligt Udgangspunkt, og det
er i Virkeligheden dette Udgangspunkt, ikke hine abstrakte
Spekulationer (resp. Sofismer), han benytter i de følgende
Dialoger, der betegner en betydningsfuld Ændring i hans
Filosofi.

I »Theaitetos« hævdes, --- i Kritik af den Paastand, at
Erkendelse skulde bestaa i passiv og momentan Sansning,
--- at al Erkendelse forudsætter en Syntese, uden hvilken
al Sondring, al Analyse vilde være umulig. De forskellige
Sansninger sidder ikke inden i os som Krigerne i den trojanske
Træhest, men de forenes i en omfattende Enhed,
hvad enten man vil kalde denne for en Sjæl eller give den
et andet Navn (p. 184 D). I selve den Kendsgerning, at
Sansefornemmelserne kan skelnes fra hverandre og komme
til at staa som lige eller forskellige, lægger der sig en
sammenfattende Funktion for Dagen (p. 185 A). Denne Betragtning
fører saa igen til (p. 185 C) at opstille Væren
(Position) og Ikkeværen (Negation), Identitet og Forskel,
samt Talbegrebet som, hvad man med et moderne Udtryk
kunde kalde fundamentale Kategorier, d. v. s. Begreber, vi
arbejder med i al vor Tænken, ligegyldigt, hvad det saa er,
vi tænker over. At der fremdrages en saadan særlig Gruppe
af Begreber eller Ideer i nøje Sammenhæng med selve
Erkendelsesvirksomheden, betegner en Fortsættelse af den
% --- p. 56 ---
\setcounter{footnote}{0}%
\opage{56}Sondring, der i »Parmenides«, mere end i tidligere Dialoger,
blev gjort mellem forskellige Grupper af Ideer.

I »Sofisten« gaar Platon nærmere ind paa selve Tankevirksomheden,
idet han gør det betydningsfulde Skridt at
undersøge Dommen som et Udtryk for Tænkning. I enhver
Dom forbindes forskellige Begreber. Dommen er en Sammenfletning
af Ideer (διὰ γὰρ τὴν ἀλλήλων τῶν εἰδῶν συμπλοκὴν
ὁ λόγος γέγονεν ἡμῖν, p. 259 E). Selve denne forbindende
Tænken, siges der videre, er et Værende, noget, som eksisterer,
og ved hvis Bortfalden al Filosofi vilde falde bort
(260 A). Muligvis ligger her et Svar paa det Spørgsmaal,
der paa det kritiske Punkt i »Parmenides« blev stillet til
den unge platoniske Sokrates: »hvad vil du nu gøre ved
Filosofien?« --- og det et Svar, der fuldt opvejer den
Række pseudodialektiske Betragtninger, som udgør anden
Del af »Parmenides«. Det var allerede i selve »Parmenides«
blevet hævdet overfor Eleaterne, at der ikke var nogen
Modsigelse i, at samme Emne kunde høre ind under forskellige
Ideer; allerede dette peger paa en vis Forbindelse
mellem Ideerne indbyrdes. Og efterat »Theaitetos« har
paavist Syntese som Forudsætning for Sansningernes indbyrdes
Forhold som lige eller forskellige, anstilles der nu
i »Sofisten«, gennem Analyse af Dommen, en Betragtning,
der viser, at de forskellige Grundbegreber eller Kategorier
ligesaa lidt som Sansningerne kan være absolut sondrede
og uafhængige af hverandre. Ikkeværen kan saaledes godt
være Prædikat ved et Værende, naar det gælder at betegne
dettes Forskel fra et andet Værende (ὄντος πρὸς ὄν ἀντίθεσις
257 E). Ligeledes kan et og samme Værende staa som ligt
og forskelligt, som hvilende eller bevæget (256 B). Til de i
»Theaitetos« nævnte Kategorier føjes nemlig her Hvile og
Bevægelse.

% --- p. 57 ---
\setcounter{footnote}{0}%
\opage{57}I sin Lære om Eros havde Platon allerede, uden at være
sig den erkendelsesteoretiske Betydning deraf bevidst, givet
et Eksempel paa saadan Sammenfletning, som enhver Dom
betyder. Thi i den stadige Stræben eller Tendens er der
baade Identitet (hvad Maalet angaar) og Forskel (formedelst
Afstanden fra Maalet), baade Hvile (i den stadige Retning)
og Bevægelse (i den urolige Higen). Og i det Begreb om
Sjæl, der opstilles i »Faidros« (p. 245 C) som det sig selv
bevægende, der derfor maa være i evig Bevægelse\footnote[1]{Begrebet αὐτοκίνητον synes at stamme fra Alkmaion. Se \textsc{Diels}: \textit{Vorsokratiker}’, p. 105.}, er der
ogsaa forudsat Hvile og Bevægelse paa én Gang. ---

Den Omstændighed, at Platon i »Sofisten« lægger en
Analyse af Domsfunktionen til Grund, minder om Kant’s
Fremgangsmaade i »Kritik der reinen Vernunft«, og moderne
Kantianere har da ogsaa ligefrem fundet kritisk
Filosofi i den sidste Periode af Platon’s Filosoferen. Nu
mødes Platon og Kant ganske vist paa dette Punkt, og det
er et Møde af verdenshistorisk Betydning. Men de mødes
derved, at Platon ender, hvor Kant begynder. Og ren Erkendelsesteoretiker
blev Platon aldrig. Erkendelsesproblemet
blev for ham ikke egentlig Grundproblemet, eller rettere,
det var for ham paa Forhaand løst i hans Metafysik.
Han fastholder trods alt sin Idelære. En Filosof er og
bliver for ham En, der gennem sin Tænkning forholder
sig til Ideen om det Værende (ὁ δὲ γε φιλόσοφος, τῇ τοῦ
ὄντος ἀεὶ διὰ λογισμῶν προσκείμενος ἰδέᾳ. Soph. 254 A).
At han med sine Analyser i »Theaitetos« og »Sofisten«
konsekvent maatte komme positivt ud over sin egen Idelære,
ligesom han negativt var kommen ud over den i
»Parmenides«, mærker han ikke. Der stillede sig ikke for
ham det Problem, der paa Grundlag af lignende Analyser
% --- p. 58 ---
\setcounter{footnote}{0}%
\opage{58}c. 2000 Aar senere stillede sig for Kant: naar vore Kategorier
tilsidst er Former for vor egen Tankevirksomhed,
hvorfra kommer saa deres Gyldighed for Genstande, som
ikke frembringes af vor Tankevirksomhed?\footnote[1]{Se \textit{Kontinuiteten i Kants filosofiske Udviklingsgang} p. 52.} Selv hvor
Platon finder nok saa megen Selvvirksomhed og nok saa
megen Uro i vor Tanke, berøres for ham Ideernes Himmel
ikke deraf. Det Spørgsmaal, hvoraf han da véd, at der
gives en saadan Himmel, opkaster han ikke.

Bonitz har allerede i sin fortræffelige Analyse af »Sofisten«
vist\footnote[2]{\textit{Platonische Studien?} p. 195 f.}, at Platon ikke har været sig bevidst, at han
var lige ved at opløse al Metafysik i Logik. Man kan finde
en første Mulighed hos ham til den betydningsfulde Vending,
der skete ved Kant’s definitive Grundlæggelse af den
kritiske Filosofi. Men denne Mulighed har Platon ikke set
og ikke benyttet, --- har hverken kunnet eller villet benytte
den. Han var stadig enig med Eleaterne i, at enhver
sand Tanke maa være en Tanke om noget, der har sin
absolute Bestaaen (Parm. p. 132 B; 135 B). ---

Dialogen »Statsmanden« har, som Navnet angiver, en
praktisk Opgave. Men for at finde den »kongelige Videnskab«
(p. 292 E) eller Kunst (p. 300 E), som maa kræves
af den rette Statsmand, foretages en omstændelig Inddeling
af Kunster og Videnskaber, en Inddeling, der tillige --- ligesom
en tilsvarende Klassifikation i »Sofisten«, --- afgiver
Eksempler paa, hvorledes Begrebers Plads i Tankens Verden
bestemmes ved deres indbyrdes Lighed og Forskel. Og
paa given Anledning fremhæves Vigtigheden af en Maalestok
for i det enkelte Tilfælde at finde det Rette. Den
praktiske Kunst eller Videnskab forudsætter ikke blot Evne
til at sammenligne Størrelser overhovedet, men særlig Evne
% --- p. 59 ---
\setcounter{footnote}{0}%
\opage{59}til at finde det bestemte Størrelsesforhold mellem forskellige
Evner og Drifter, som maa være tilstede, hvis de Opgaver,
Etik og Statslære stiller, skal kunne ske Fyldest.
Det gælder for alle praktiske Bestræbelser, at de kun kan
føre til, hvad der er skønt og godt, naar det rette Maal
overholdes (τὸ μέτριον σώζουσαι πάντ’ ἀγαθὰ καὶ καλὰ ἀπεργάζονται.
 Polit. p. 284 B). --- Maalekunstens etiske Betydning
var allerede fremhævet i »Protagoras« (p. 356—357)
og i »Gorgias« (p. 507—508).

Denne sidste Betragtning føres videre i »Filebos«, der,
i Forbindelse med »Timaios«, giver den erkendelsesteoretiske
Afslutning af den Tankebevægelse, der begyndte i
»Parmenides«, og derved staar som det afgørende Udtryk
for Platon’s senere Filosofi. Nogle Forfattere ser i denne
Dialog den »Filosofos«, som Platon i »Sofisten« (p. 217 A)
havde bebudet som tredie Led i en Trilogi, hvis to første
Led skulde være »Sofisten« og »Statsmanden«.

I »Sofisten« blev en indre Sammenhæng mellem Ideerne
godtgjort ved, at al Tankevirksomheds Grundform, Dommen,
er en Forbindelse af Begreber, en »Blanding«. I
»Statsmanden« fremhævedes Nødvendigheden af et bestemt
Størrelsesforhold, et bestemt Maalforhold, naar det gælder Afgørelsen
af praktisk Værdi. I »Filebos« vises det nu, hvorledes
det ved al Forstaaelse og ved al Vurdering gælder at
finde et bestemt Forhold mellem en Grænse (Regel) og et
Ubegrænset, som skal begrænses. Ved det Ubegrænsede
(τὸ ἄπειρον) tænkes særlig paa de usammenhængende Forskelligheder
(smlg. Polit. p. 273 D: ὁ τῆς ἀνομοιότητος τόπος
ἄπειρος ὤν). Platon’s Kamp for sine Ideer er stadig en
Kamp mod kaotiske Forskelsrækker. Hverken i Ideernes
eller Emnernes Verden kan saadanne Forskelsrækker taales.
Alt hvad der er, bestaar af Enhed og Mangfoldighed og
% --- p. 60 ---
\setcounter{footnote}{0}%
\opage{60}har Grænse og Ubegrænsethed i sin Natur. (Fileb. p. 16 C).
Grænsen lægger sig for Dagen derved, at Identitet, Tal og
Maal raader (25 A). Dialogens egentlige Emne er psykologisk
= etisk, derfor spiller Maalestokken for Vurdering en
større Rolle end Maalestokken for Forstaaelse, skønt der
er fælles Forudsætninger for begge. Medens Matematiken i
»Staten« opfattes som et Trin af Erkendelse, der hviler
paa visse Forudsætninger, hvis Grund først Dialektiken
afslører, afgiver den i »Filebos« Middel og Form for Bestemmelse
af Tingenes Væsen, ganske særlig for Bestemmelse
af Værdier (»det Gode«). Begrebet om det rette Maal
(τὸ μέτριον) bliver nu et centralt Begreb. Derved, at »det
Gode« forudsætter et maalbestemt og harmonisk Forhold
(μετριότης καὶ συμμετρία), fremtræder det som skønt (p. 64 E).

Platon kommer her uvilkaarlig til at antyde Sammenhængen
mellem Begrebet Værdi og Begrebet Totalitet. Thi
spørges der, hvad da Følgen bliver af, at det rette Maal
ikke fastholdes, svares der, at det fører til Opløsning i et
ulykkeligt Kaos (p. 64 DE smlgn. Polit. p. 273 D). I vort
moderne Sprog kan vi udtrykke dette ved, at enhver Helhed
kræver visse indre og ydre Betingelser for sin Bestaaen,
og at den Vurdering, vi anstiller af personlige og
sociale Helheders Tilstande og Handlinger, tilsidst begrundes
ved disse Tilstandes og Handlingers Indgriben i
Betingelserne for vedkommende Helheds Bestaaen og Udvikling\footnote[1]{Smlgn. \textit{Totalitet som Kategori,} p. 66---73.}.
Muligheden af en rationel Etik beror paa Muligheden
af nøjagtig Bestemmelse af saadanne Livsbetingelser.
I »Staten« har Platon søgt at vise, at den enkelte Personlighed
og Statens Livsbetingelser staar i harmonisk Forhold
til hverandre, naar begge Helheder naar deres fuldkomne
Udvikling. I sine senere Aar synes han at have lagt hele
% --- p. 61 ---
\setcounter{footnote}{0}%
\opage{61}Vægten paa de kvantitative Forhold, der maa fyldestgøres,
for at disse Livsbetingelser kan være tilstede. Det berettes,
at de, der kom for at høre Platon’s Forelæsning over »det
Godes« Begreb blev i høj Grad skuffede ved kun at høre
tale om Aritmetik og Astronomi, om »Grænsen« og om »det
Ene«. Paa den anden Side har han i sine ovenfor anførte
Udtalelser i »Statens« syvende Bog og i det syvende Brev
hævet »det Godes Ide« saa højt over al Viden, at han
derved berøver sig det rationelle Grundlag for Etiken og
tillige sønderriver det Baand mellem det Gode og det
Skønne, der staar fast for ham, saalænge han færdes indenfor
menneskelig Erfaring. Det mystiske og det rationelle
Element i hans Filosofi træder her i skarp Modsætning til
hinanden. ---

Ligesaa lidt som den erkendelsesteoretiske Tankegang i
»Sofisten« fører Læren om Begrænsning og om Nødvendighed
af Maalbestemthed for enhver Helhed (»Blanding«)
Platon til at opgive Idelæren. Den højeste Viden er stadig
den om »det i Sandhed Værende, det, der stedse forholder
sig paa samme Vis« (p. 58 A). Idelæren ligger stadig i Baggrunden,
efterat det har vist sig nødvendigt at anerkende
den Forbindelse, »Sofisten« har paavist gennem Analyse
af Dommen, og den Blandingskarakter (Helhedskvalitet),
som »Filebos« har paavist ved alt, der bestaar og har Værdi. ---

I »Timaios« er det en Hovedtanke, at Forbindelsen af
Identitet og Forskel er udtrykt i Universets Bygning. I
Himmelens Ækvator fremtræder Identiteten gennem Himmelkuglens
ensformige Bevægelse; i Ekliptika fremtræder Forskellighederne
gennem Solens og Planeternes Bevægelser.
Svarende hertil lægger Identiteten sig for Dagen i Menneskets
Tænkning, Forskelligheden i Sansningen. Den samme
Modsætning gør sig gældende i den Universet omspændende
% --- p. 62 ---
\setcounter{footnote}{0}%
\opage{62}Verdenssjæl, der betinger Naturens matematiske Lovmæssighed.
Denne Forbindelse af Identitet og Forskel svarer til
»Sofistens« Forbindelse af Begreber i Dommen og til »Filebos’s«
Forbindelse af Grænse og Ubegrænset. Det var den
i de nævnte Dialoger fremstillede nye Lære, der gjorde det
muligt for Platon at faa en Ramme, indenfor hvilken han
kunde ordne det fysiske, astronomiske og biologiske Stof,
han i sine senere Aar samlede. Man har med Rette bemærket,
at det vilde være en Misforstaaelse at tro, at der
i Akademiet kun dyrkedes Filosofi, Matematik og Astronomi.
Ogsaa biologiske Studier optog Sindene, og maaske
lagdes allerede her Grunden til Aristoteles’ Interesse for
% PRINTED AS IS: „BURNÉT“ with an accent (note 1), p. 62, seen at 800 dpi
Biologi\footnote[1]{Se \textsc{Burnét}: \textit{Greek Philosophy. I. Thales to Platon,} p. 223.}. Medens han i sin klassiske Idelære (i »Symposion«,
»Faidros«, »Staten«) lod sig nøje med at stige op
fra Emner til Ideer og gaa fra Ide til Ide indenfor Ideernes
Verden, havde han nu fremsat Tanker, der gjorde det
muligt for ham at tillægge Emnerne anden Betydning end
den blot at være Springbrædt ved Farten op til Ideverdenen,
og det blev en Opgave for ham at finde en Sammenhæng
mellem Emnerne indbyrdes. Paa Grund af de ufuldkomne
Hjælpemidler maatte et Forsøg i denne Retning
komme til at tage sig barnlig ud, og Platon gør rigelig
Brug af mytologiske Forestillinger.

Baade fra filosofisk og fra filologisk Side har man paavist,
at der i »Timaios« gør sig en Indflydelse af Demokrits
Tanker gældende, særlig af hans Atomlære. Den
strenge Lovmæssighed, som Platon hævder (skønt han
underordner den under et teleologisk Synspunkt), gælder
ikke blot for de sanselige Emners Verden, men ogsaa for de
i sig usanselige Elementer, den tilsidst bestaar af. Allerede
Demokrit’s Atomer var usanselige, og man var allerede i
% --- p. 63 ---
\setcounter{footnote}{0}%
\opage{63}Oldtiden opmærksom paa den Lighed mellem Platon og
Demokrit, at de begge, hver paa sin Vis, forklarede de
sanselige Kvaliteter ved usanselige Aarsager. Platon gaar
længere tilbage end Demokrit, idet Atomerne efter ham
maa antages at være sammensatte af Elementer, der kun
kan karakteriseres rent matematisk, og til Grund for disse
Elementer ligger igen (som en Slags Atomer af tredie Grad)
finere Elementer, der »kun kendes af Gud og af dem han
elsker« (p. 53 D). Paa interessant Maade er her Atombegrebets
Relativitet antydet i mytologisk Form\footnote[1]{\textsc{Léon Robin} (\textit{La Place de la Physique dans la Philosophie de Platon,} Paris 1919), som stærkt har henledet Opmærksomheden paa det anførte Sted i »Timaios«, søger at vise, at der ved Atomerne af tredie Orden maa tænkes paa de Idealtal og Idealfigurer, hvoraf Platon (ifølge Aristoteles) saa meget optages i sin sidste, pytagoræiserende Periode. --- Gennem Platon’s Naturfilosofi gaar, ved Siden af det meget Fantastiske, en Stræben efter at føre alle Kvaliteter tilbage til rent kvantitative Forhold --- skønt ogsaa dette sker paa fantastisk Maade.}.

Idet Platon saaledes indfører en tredie Verden i sin
Filosofi ved Siden af Ideernes Verden og de sanselige Emners
Verden, foretager han i Overensstemmelse dermed en
mærkelig Ændring i sin Erkendelsesteori. Atomerne kan
hverken tænkes (paa den Maade, paa hvilken Ideerne kan
tænkes) eller sanses (som Emnerne med deres Kvaliteter).
Kun Analogiens Vej staar da aaben. Det er den, Platon
peger paa, naar han siger, at det, der ligger til Grund for
de sanselige Kvaliteter, maa opfattes »uden Sansernes Hjælp
ved en Slags uægte Tænkning« (p. 52 B). (Endnu p. 27 D
nævnes kun Tænkning og Sansning uden Antydning af en
tredie Art Viden)\footnote[2]{Se min Afhandling \textit{Platon og Demokritos} (Mindre Arbejder. Tredie Række), p. 99.}.

Det er dog ikke sikkert, at den Trehed af Ideer, sanselige
Emner og Atomer af forskellige Grader, som har
% --- p. 64 ---
\setcounter{footnote}{0}%
\opage{64}fremkaldt en tilsvarende Trehed indenfor Erkendelsesteorien,
blot skyldes Indflydelse af Demokrit’s Lære. Léon
Robin, som især støtter sig til »Timaios« p. 52 C, har (anf.
Skrift p. 26) med Rette fremhævet, at naar Emnerne ikke
kunde udledes af Ideerne, maatte der forudsættes en Faktor,
som kunde gøre det muligt, at Emnerne trods deres
Lighed med Ideerne dog kunde have Egenskaber, hvor de
afveg fra dem. Der maa være noget, hvorved Afbilledet er
forskelligt fra Forbilledet; --- der maa altsaa gives noget
fra Ideerne forskelligt (ἕτερον τι p. 52 C). Hvis der ikke
gaves andet end Ideerne, vilde deres Afbilleder være uforklarlige.
Det er ved et Postulat, der begrundes ved Forskellen
mellem Ideer og Emner, at Platon ud fra sin tidligere
Lære kommer til Principerne for sin Naturfilosofi i
»Timaios«. Han drager egentlig her Konsekvensen af den
tredie Indvending i »Parmenides« mod Idelæren, --- et
Vidnesbyrd om, at han har følt sig ramt af den. ---

Er der i den Tankegang, Platon kommer ind paa i
»Timaios«, Træk, der gør det forstaaeligt, at Renæssancens
Naturforskere, især Galilei, kunde finde en Aandsfrænde i
ham, saa forstaas dette endnu klarere ved den Problemstilling,
Platon ifølge Simplikios skal have anvist Astronomien.
Han finder nemlig dennes Opgave i at udforske,
hvilke regelmæssige Bevægelser man maa antage i Universet,
for at de i Erfaringen foreliggende Fænomener kan
komme til deres Ret. Det er en Maade at spørge paa, som
vistnok skriver sig fra Pytagoræerne, men Platon gav den
en strengere Form, og Platon’s Ven Eudoxos søgte at løse
den Opgave, den stillede. Som Duhem, der særlig har
undersøgt dette Punkt, siger: »De fysiske Videnskabers
Metode er bleven defineret af Platon og af Pytagoræerne
paa hans Tid med en Tydelighed og en Nøjagtighed, der
% --- p. 65 ---
\setcounter{footnote}{0}%
\opage{65}ikke er blevne overgaaede«\footnote[1]{\textit{Le Système du Monde. I. De Platon à Copernik.} 1913, p. 129.}. Hvis dette er rigtigt, har Platon
gennem sine erkendelsesteoretiske Fordringer til Videnskaben
øvet en indirekte Indflydelse paa Naturvidenskabens
Udvikling, ligesom han ifølge Zeuthen har øvet en saadan
Indflydelse paa Matematikens Udvikling. ---

Det er et stort Spørgsmaal, om Platon vilde være kommen
ind paa de betydningsfulde erkendelsesteoretiske Ideer,
han har udviklet i sine senere Dialoger, Ideer, der giver
disse Dialoger en stor Interesse for Tænkningens Historie
trods deres kunstneriske Ufuldkommenhed, hvis han ikke
med saa stor Styrke og Sandhedskærlighed havde gennemtænkt
de Indvendinger, der kunde rejses, og tildels er blevne
rejste mod hans Idelære. De førte ham ikke til at opgive
denne, men de har virket som en Spore til fornyet Tankearbejde
og har derigennem virket til at give hans Tænkning
en større Universalitet, end den vilde have haft, naar
den dogmatisk havde afsluttet sig om den oprindelige Idelære
som fuldstændigt Udtryk for Sandheden.

% --- Indhold (PDF 68, p. [66]), read at the image (m1 F68, reader S).  The Indhold gives the Indledning as p. 1 (the text
% begins on p. 3) and chapter 3's title as „Begrebet det Pludselige i Platon’s Filosofi“ (the head on p. 38: „... (τὸ ἐξαίφνης) i
% Platons Filosofi.“); both are recorded as printed.
\clearpage
\begin{center}{\large I N D H O L D}\\[0.3em]\rule{3em}{0.4pt}\end{center}
\hfill{\small Side}\par
\tocline{}{Indledning}{1}
\tocline{1.}{Kritik af Idelæren}{6}
\tocline{2.}{Dialogens anden Del}{26}
\tocline{3.}{Begrebet det Pludselige i Platon’s Filosofi}{38}
\tocline{4.}{Senere Dialoger i Forhold til Problemerne i Parmenides}{53}
\begin{center}\rule{3em}{0.4pt}\\[-0.8em]\rule{3em}{0.4pt}\end{center}
\vfill
\begin{center}{\small Forelagt paa Mødet den 16. Januar 1920.\\ Færdig fra Trykkeriet den 15. Juni 1920.}\end{center}
\end{document}
